\documentclass[11pt]{article}

\usepackage{geometry}
\geometry{margin=1in}

\usepackage{amsmath,amssymb,amsfonts,amsthm}
\usepackage{mathtools}
\usepackage{hyperref}
\usepackage{graphicx}
\usepackage{float}
\usepackage{booktabs}
\usepackage{tikz-cd}
\usepackage{cleveref}
\usepackage{subcaption}
\usepackage{placeins}
\usepackage{enumitem}
\usepackage{array}
% xr provides \externaldocument so cross-document \ref{S-...} works.
% Plain xr (not xr-hyper) is used because xr-hyper imports labels in a
% hyperref-extended format that cleveref's templabel parser does not accept.
% The cost is that cross-document references show as text only rather than
% as live hyperlinks; the in-document \ref number resolution still works
% and there are no undefined-reference errors.
\usepackage{xr}

\hypersetup{colorlinks=true, linkcolor=blue, citecolor=blue, urlcolor=blue}

%---------------------------
% Theorem environments
%---------------------------
\numberwithin{equation}{section}

\newtheorem{theorem}{Theorem}[section]
\newtheorem{lemma}[theorem]{Lemma}
\newtheorem{proposition}[theorem]{Proposition}
\newtheorem{corollary}[theorem]{Corollary}
\theoremstyle{definition}
\newtheorem{definition}[theorem]{Definition}
\newtheorem{assumption}[theorem]{Assumption}
\theoremstyle{remark}
\newtheorem{remark}[theorem]{Remark}

\newcommand{\proofsuppref}[1]{%
  \par\noindent\emph{Proof.} See the online supplement, Section~\ref{S-#1}.\par\medskip
}
\newcommand{\proofappendixref}[1]{\proofsuppref{#1}}


%---------------------------
% Notation shortcuts
%---------------------------
\newcommand{\K}{\mathcal{K}}
\newcommand{\Hil}{\mathcal{H}}
\newcommand{\bd}{\partial}
\newcommand{\Id}{\mathbf{1}}
\newcommand{\Proj}{\mathbf{P}}
\newcommand{\spec}{\mathrm{spec}}
\newcommand{\dist}{\mathrm{dist}}
\newcommand{\rank}{\mathrm{rank}}
\newcommand{\tr}{\mathrm{tr}}
\newcommand{\Ran}{\mathrm{Ran}}
\newcommand{\opnorm}[1]{\left\lVert#1\right\rVert}
\newcommand{\e}{\varepsilon}
\newcommand{\sgn}{\mathrm{sgn}}
\newcommand{\Indint}{\operatorname{Ind}_{\mathrm{int}}}
\newcommand{\Ch}{\operatorname{Ch}}
\newcommand{\hexagon}{\mathrm{hex}}

% Resolve cross-document \ref{S-...} labels via xr.
\externaldocument[S-]{paper_supplement}

%---------------------------
% Example name macros
%---------------------------
\newcommand{\ExOneName}{Disk-like patch (1 boundary)}
\newcommand{\ExTwoName}{Annulus-like band (2 boundaries)}
\newcommand{\ExThreeName}{Three-boundary patch (disk with two holes)}

%---------------------------
% Figure safety helper
%---------------------------
\newcommand{\safeincludegraphics}[2][]{%
  \IfFileExists{#2}{\includegraphics[#1]{#2}}{%
    \fbox{\begin{minipage}[c][0.18\textheight][c]{0.92\linewidth}
      \centering Missing figure file\\[0.3em]\texttt{\detokenize{#2}}
    \end{minipage}}%
  }%
}

\title{Boundary and Interface Classes for Universal Chern Hamiltonians\\
via Doubling and Jamming}

\author{
Yitzchak Shmalo\thanks{Corresponding author: \texttt{yitzchak.shmalo@gmail.com}.}\\[0.6em]
Einstein Institute of Mathematics\\
Hebrew University of Jerusalem
}
\date{\today}

\begin{document}
\maketitle

\begin{abstract}
Higson and Prodan's universal Chern Hamiltonians are closed-surface models built from the
simplicial boundary operator and Poincar\'e duality. We treat the boundary and interface classes
obtained from finite surfaces with boundary by doubling the surface and jamming the reflected copy.
For every finite triangulated oriented surface with boundary we construct an intrinsic
Poincar\'e--Lefschetz Hamiltonian, compare it with the original-side compression of the doubled
closed Hamiltonian, identify the discrepancy as boundary-collar supported, and prove
Schur-complement estimates for finite-\(M\) jamming. Along families with negligible collar
dimension these estimates identify the normalized bulk spectral limits of the intrinsic,
compressed, and jammed models. The headline Roe-algebra theorem identifies the doubled-and-jammed
interface class in \(K_1(C^*(Y\subset X))\) with
\(\partial_Y(E_+(\pi(P_\pm)))\), using the Higson--Prodan closed universal-model theorem
\cite{HigsonProdanUniversalPRL,HigsonProdanUniversalArxiv} for the closed-bulk gap and Chern
value; this \(K_1\)-class identity is unconditional modulo HP. The integer interface winding
\(\pm1\) is proved internally in the periodic triangular Toeplitz case via the
\(K_1\)-to-\(\mathbb Z\) Toeplitz extension, and is stated as a conditional corollary for
general partitions, conditional on the existence of an integer-valued Chern-jump pairing for
the interface ideal. The periodic triangular Bloch class is a fully internal benchmark: its
closed-bulk gap, Chern signs, and straight Toeplitz winding are proved in the paper.

\medskip\noindent
\textbf{Dependency categorization (sharp framing).} The results split into three buckets
according to which external theorems each one uses. \emph{(A) Fully internal:} the
finite-dimensional boundary theory and Schur-complement estimates, the internal bulk spectral
equivalence, and the periodic triangular benchmark (gap, Chern signs, and integer Toeplitz
winding for the straight periodic interface). \emph{(B) Internal modulo Higson--Prodan:} the
\(K_1\)-class identity \(\Indint=\partial_Y(E_+(\pi(P_\pm)))\) for general HP-covered boundary
families. \emph{(C) Conditional on an external integer pairing:} the integer evaluation
\(\pm 1\) of this \(K_1\)-class outside the periodic Toeplitz case. The Higson--Prodan
closed-bulk gap-and-Chern theorem is the \emph{only} external theorem invoked anywhere in this
paper; the Ewert--Meyer language is used to fix Mayer--Vietoris terminology only, and the
required compatibility with the Roe connecting map is proved here. For flat-torus closed
doubles, both the closed-bulk gap (\Cref{S-thm:internal_torus_gap}) and the closed-bulk Chern
integer (\Cref{S-thm:internal_torus_chern}) are proved internally; both two-band benchmarks
on the same flat-torus quotient (hexagonal and square) are also proved internally for both the
gap input (\Cref{S-thm:internal_hex_gap} and \Cref{S-thm:internal_square_gap}) and the Chern
integer input (\Cref{S-thm:internal_hex_chern} and \Cref{S-thm:internal_square_chern}); the
three internal Chern theorems are unified by a single general lattice-Chern lemma
(\Cref{S-thm:internal_bc_chern_general}) for any continuously parametrized self-adjoint Bloch
family on \(\mathbb T^2\) with uniform Fermi gap, with the three benchmark-specific theorems
recovered as corollaries
(\Cref{S-cor:bc_chern_torus_corollary,S-cor:bc_chern_hex_corollary,S-cor:bc_chern_square_corollary}).
A standard internal locality tool for the bounded-geometry case, the Combes--Thomas
exponential off-diagonal decay of the Fermi projection on the closed universal Hamiltonian
\(H_\pm(K)=(B+B^\dagger)\pm S_K\), is also recorded internally
(\Cref{S-thm:internal_combes_thomas}); it converts a uniform Fermi gap on a bounded-geometry
triangulation \(K\) into a uniform exponential decay rate for \(P_\pm(K)\), but does not by
itself supply the bounded-geometry gap. The general bounded-geometry case still cites HP for
the gap input.

\medskip
\noindent Companion
numerical code and the online supplement are available at the project repository
(\url{https://github.com/yspennstate/Doubling-and-Jamming-A-Universal-Chern-Hamiltonian-on-Triangulated-Surfaces-with-Boundary});
additional hexagonal and square two-band symbols, robustness variants, full appendix proofs,
and the numerical record are collected in that online supplement.
\end{abstract}

%====================================================================
\section{Introduction}
%====================================================================

\subsection{Universal Chern Hamiltonians from simplicial data}
The Haldane model is the prototypical two-dimensional class-A topological insulator
\cite{Haldane1988}; for background on the tenfold-way classification and Chern phases, see
\cite{SchnyderRyuFurusakiLudwig,KitaevPeriodic,QiZhangRMP,ProdanSchulzBaldesBook}. On irregular
discretizations, direct magnetic-field or Peierls constructions can lead to mobility-gap or
defect-state issues that complicate the extraction of clean topological gaps; see, for example,
\cite{Mitchell2018,BourneProdan2018,BourneMesland2019,ProdanSchulzBaldesBook}. Rigorous
bulk--edge and bulk--interface correspondences in quantum Hall and related class-A systems are
treated in
\cite{ElbauGraf2002,ElgartGrafSchenker2005,GrafPorta2013,BourneKellendonkRennie,KubotaControlled,EwertMeyer2019,AlldridgeMaxZirnbauer2020,LudewigThiangEdge}.

Higson and Prodan introduced a different route
\cite{HigsonProdanUniversalPRL,HigsonProdanUniversalArxiv}. For a triangulation of a \emph{closed}
oriented surface, one builds a local tight-binding Hamiltonian from two pieces of
simplicial-topological data: the simplicial boundary operator $B$ and a Hermitian Poincar\'e-duality
operator $S$ constructed from cap product with the fundamental cycle. The resulting Hamiltonians
\begin{equation}\label{eq:Hpm_intro}
H_\pm=(B+B^\dagger)\pm S
\end{equation}
are local and, in the sequential-refinement limit, define nontrivial topological bulk classes with
Chern numbers $\pm1$ \cite{HigsonProdanUniversalPRL,HigsonProdanUniversalArxiv}. The underlying
Hilbert--Poincar\'e framework uses the algebraic identity \(BS+SB^\dagger=0\), rooted in analytic
surgery theory and the work of Higson and Roe
\cite{HigsonRoeI,HigsonRoeII,HigsonRoeIII,WallSurgery}.

\subsection{The boundary case}
If the input triangulation $\K$ has boundary, there is no closed fundamental $2$-cycle supported
on the faces of $\K$, so the closed-case cap-product recipe for $S$ is not directly available.
Compact manifolds with boundary instead carry a relative fundamental class and satisfy
Poincar\'e--Lefschetz duality \cite{Hatcher}. We use this relative class to define an intrinsic
boundary Hamiltonian on $\K$, and then compare it with a second construction that reduces the
boundary problem to the closed case by simplicial doubling and jamming the reflected side.

\subsection{Doubling and jamming}
Given a triangulated surface-with-boundary $\K$, let $\overline\K$ be an orientation-reversed copy.
Gluing $\K$ and $\overline\K$ along $\bd\K$ produces the closed double $D(\K)$. We then build the
closed universal Hamiltonians on $D(\K)$ and suppress the reflected copy by adding a large on-site
potential on the copy degrees of freedom. This is the \emph{doubling-and-jamming} construction. The
finite-\(M\) analysis gives error bounds for this approximation. For refinement families whose
closed doubles satisfy the hypotheses of the Higson--Prodan closed-bulk theorem
(\Cref{thm:hp_imported}; hypothesis check for the doubled families in
\Cref{lem:hp_hypothesis_check}), that theorem supplies the closed-bulk gap and Chern class.
Section~\ref{sec:pathb_internalization} proves, for the original/copy partition, that the Roe
connecting map used to define the interface class is the Ewert--Meyer bulk--interface map for
the doubled-jammed data.

\subsection{Theorem map}
Let
\[
H_{\pm,M}=H_\pm(D(\K))+M\Proj_{\mathrm{copy}}
\]
be the jammed Hamiltonian, written in block form as
\[
H_{\pm,M}=
\begin{pmatrix}
A&C\\
C^\dagger&D+M\Id
\end{pmatrix}
\quad\text{on}\quad
\Hil_{\mathrm{orig}}\oplus\Hil_{\mathrm{copy}}.
\]
\begin{enumerate}[leftmargin=2em]
\item The relative fundamental class of \((\K,\bd\K)\) gives an intrinsic
Poincar\'e--Lefschetz boundary Hamiltonian on \(\K\). The simplicial double gives a second
Hamiltonian by closed Higson--Prodan theory; its original-side compression agrees with the
relative Hamiltonian up to a boundary-collar operator, and the jammed operator converges to this
compression at low energy.
\item The finite-\(M\) comparison is made quantitative by Schur-complement estimates:
compressed resolvents, full block resolvents, first-order effective operators, bounded-energy
spectral localization, leakage, Morse-index stability, and contour spectral projections. Combined
with the collar-rank comparison, these estimates give a fully internal bulk spectral-measure
equivalence between the intrinsic Poincar\'e--Lefschetz model, the doubled compression, and the
finite-\(M\) jammed model along any family with negligible boundary-collar dimension and
\(M\to\infty\).
\item For refinement families whose closed doubles satisfy the hypotheses of
\Cref{thm:hp_imported} (verified in \Cref{lem:hp_hypothesis_check}), the
HP gap gives the quotient by the interface ideal and hence the global Roe-algebra interface class
\[
\Indint(J_{\pm,M}^\infty;T^+,T^-)
=
\partial_Y\!\bigl(E_+(\pi(P_\pm))\bigr)
\in K_1(C^*(Y\subset X)).
\]
Theorem~\ref{thm:pathb_roe_winding}, proved in Section~\ref{sec:pathb_internalization}, identifies
this class as \(\partial_Y(E_+(\pi(P_\pm)))\). In the straight periodic Toeplitz case the
integer pairing is internal (Corollary~\ref{cor:toeplitz_winding}) and gives winding
\(\pm1\); for general partitions the integer pairing is recorded as a conditional corollary
(Corollary~\ref{cor:conditional_winding}).
\item For the periodic triangular-lattice class, the closed-bulk gap and Chern input are checked
directly. We write down the \(6\times6\) Bloch symbol, prove a uniform zero gap from an exact
determinant formula, compute the Chern signs from the isolated massless crossing, and evaluate
primitive straight periodic interfaces by the Toeplitz extension. Finite covers, annular families,
finite-corner quotients, finite-defect variants, and tail-split component maps are treated in the
online supplement, Section~\ref{S-app:section5_variants}.
\end{enumerate}

\paragraph{Logical dependencies.}
The proof uses the following implications; no converse is assumed.
\[
\begin{array}{ll}
\mathrm{(I)} & \text{internal finite-dimensional results: the relative boundary operator,}\\
& \text{the collar comparison, and the large-\(M\) Schur-complement estimates;}\\[0.2em]
\mathrm{(HP)} & \text{the Higson--Prodan closed-bulk gap and Chern theorem for the doubles;}\\[0.2em]
\mathrm{(EM)} & \text{the internal identification of the original/copy Roe boundary map}\\
& \text{with the Ewert--Meyer pairing (Theorem~\ref{thm:pathb_roe_winding}).}
\end{array}
\]
From (I) and the gap part of (HP), the Roe construction supplies the global interface
\(K_1\)-class. From the Chern part of (HP) and (EM), Theorem~\ref{thm:pathb_roe_winding}
identifies the \(K_1\)-class as \(\partial_Y(E_+(\pi(P_\pm)))\). Thus the headline conclusion
of Theorem~\ref{thm:pathb_roe_winding} is the unconditional
\(K_1\)-class identity
\(\mathrm{(I)}+\mathrm{(HP)}+\mathrm{(EM)}\Longrightarrow
\Indint=\partial_Y(E_+(\pi(P_\pm)))\). The integer evaluation
\(\operatorname{Wind}=\pm1\) is internal in the periodic triangular Toeplitz case
(Corollary~\ref{cor:toeplitz_winding}) and conditional in general
(Corollary~\ref{cor:conditional_winding}). The ingredients (I) and (EM) are proved here. The
ingredient (HP) is exactly the cited closed theorem of Higson--Prodan. In the periodic
triangular section, the gap and Chern signs are proved directly, and primitive straight
periodic interfaces are evaluated internally by the Toeplitz extension.

\begin{table}[H]
\centering
\small
\caption{Theorem map. ``Internal'' means proved here from the inputs shown; ``modulo HP'' means
that the only cited closed-bulk theorem used is Higson--Prodan's published PRL/arXiv theorem.}
\label{tab:main_theorem_map}\label{tab:hypothesis_ledger}
\begin{tabular}{>{\raggedright\arraybackslash}p{0.30\linewidth}
                >{\raggedright\arraybackslash}p{0.33\linewidth}
                >{\raggedright\arraybackslash}p{0.26\linewidth}}
\toprule
Category and conclusion & Inputs & Status \\
\midrule
Finite-dimensional boundary theory: relative boundary Hamiltonian and collar comparison & Finite
triangulated surface with full
boundary subcomplex & Internal, \Cref{thm:relative_boundary_comparison} \\
Finite-dimensional jamming theory: finite-\(M\) decoupling, leakage, contour projections, Morse
index & Finite block separation
hypotheses & Internal, Section~\ref{sec:decoupling} \\
Internal bulk spectral measures: PL model, doubled compression, and finite-\(M\) jammed model have
the same normalized bulk spectral limits & Uniform finite-dimensional bounds, \(M\to\infty\), and
boundary-collar dimension \(o(\dim\Hil(\K_n))\) & Internal,
\Cref{thm:internal_bulk_spectral_equivalence} \\
Roe-algebra \(K_1\)-class identity: global interface
\(\partial_Y(E_+(\pi(P_\pm)))\in K_1(C^*(Y\subset X))\) & The Higson--Prodan closed-bulk
gap/Chern theorem and the internal original/copy Roe/Ewert--Meyer identification & Internal
modulo HP, \Cref{thm:pathb_roe_winding} \\
Toeplitz integer winding \(\pm1\) for straight periodic interfaces & Periodic triangular gap,
Chern signs, BDF--Toeplitz extension & Internal,
\Cref{cor:toeplitz_winding,thm:periodic_toeplitz_pairing} \\
Integer winding for general partitions & Hypothesized Chern-jump-compatible integer pairing on
\(K_1(\mathcal I)\) & Conditional on the pairing's existence,
\Cref{cor:conditional_winding} \\
Internal periodic examples: triangular closed-bulk input & Periodic triangular Bloch symbol &
Internal gap and Chern
signs, \Cref{thm:periodic_triangular_bulk,thm:periodic_triangular_chern} \\
Flat-torus closed-bulk gap (uniform triangulation \(\K_N^\triangle\) on
\(T_N=\mathbb R^2/(N\mathbb Z)^2\)) & Bloch decomposition of finite torus + periodic triangular
gap & Internal, \Cref{S-thm:internal_torus_gap}; replaces HP gap input on the flat-torus
instance \\
Flat-torus closed-bulk Chern integer
\(\Ch_N(P_\pm(\K_N^\triangle))=\mp1\) (\(N\ge N_0\)) &
Sampled-Bloch FHS lattice formula + periodic triangular continuous Chern + FHS convergence
(\cite{FukuiHatsugaiSuzuki,BellissardVanElstSchulzBaldes1994}) & Internal,
\Cref{S-thm:internal_torus_chern}; replaces HP Chern input on the flat-torus instance \\
Flat-torus hexagonal closed-bulk gap (hexagonal benchmark Bloch model
\(H_{\hexagon}(k;m)\) on \(T_N=\mathbb R^2/(N\mathbb Z)^2\), \(m\ne 0\)) & Bloch decomposition
of finite torus + symbol-level hexagonal determinant lower bound (\(2\)-band) & Internal,
\Cref{S-thm:internal_hex_gap}; replaces HP gap input on the hexagonal flat-torus
instance \\
Flat-torus square closed-bulk gap (square benchmark Bloch model
\(H_{\square}(k;m)\) on \(T_N=\mathbb R^2/(N\mathbb Z)^2\), \(m\ne 0\)) & Bloch decomposition
of finite torus + symbol-level square determinant lower bound (\(2\)-band) & Internal,
\Cref{S-thm:internal_square_gap}; replaces HP gap input on the square flat-torus
instance \\
Flat-torus hexagonal closed-bulk Chern integer
\(\Ch_N(P_-(\K_N^{\hexagon};m))=-\sgn(m)\) (\(N\ge N_0^{\hexagon}(m)\)) &
Sampled-Bloch FHS lattice formula + symbol-level hexagonal continuous Chern + FHS convergence
(\cite{FukuiHatsugaiSuzuki,BellissardVanElstSchulzBaldes1994}) & Internal,
\Cref{S-thm:internal_hex_chern}; replaces HP Chern input on the hexagonal flat-torus
instance \\
Flat-torus square closed-bulk Chern integer
\(\Ch_N(P_-(\K_N^{\square};m))=+\sgn(m)\) (\(N\ge N_0^{\square}(m)\)) &
Sampled-Bloch FHS lattice formula + symbol-level square continuous Chern + FHS convergence
(\cite{FukuiHatsugaiSuzuki,BellissardVanElstSchulzBaldes1994}) & Internal,
\Cref{S-thm:internal_square_chern}; replaces HP Chern input on the square flat-torus
instance \\
General lattice-Chern convergence (any \(C^4\) self-adjoint Bloch family on
\(\mathbb T^2\) with uniform Fermi gap and constant negative rank;
\(\Ch_N(P)=\Ch(P)\) for \(N\ge N_0(\gamma,\Lambda,d)\)) & Uniform gap, uniform norm bound,
constant negative rank, \(C^4\) smoothness; FHS plaquette identity; Kato analytic perturbation
(\cite{FukuiHatsugaiSuzuki,BellissardVanElstSchulzBaldes1994,KatoPerturbation}) & Internal,
\Cref{S-thm:internal_bc_chern_general}; unifies the three flat-torus Chern theorems above as
corollaries (\Cref{S-cor:bc_chern_torus_corollary,S-cor:bc_chern_hex_corollary,S-cor:bc_chern_square_corollary}) \\
Internal Combes--Thomas exponential decay for the universal Hamiltonian on a bounded-geometry
triangulation (\(|\langle\delta_v,P_\pm(K)\delta_w\rangle|\le C(\alpha)e^{-\alpha d(v,w)}\)
for \(\alpha<\alpha_0(\gamma,\Lambda,R,\Delta)\)) & Uniform Fermi gap
\(\dist(0,\spec H_\pm(K))\ge\gamma\), uniform hopping range \(R\), uniform norm bound \(\Lambda\),
bounded-geometry parameters \((\Delta,L,R_0)\); imaginary-shift conjugation, Neumann series for
the resolvent, Riesz contour integral
(\cite{CombesThomas1973,AizenmanWarzel2015,ProdanSchulzBaldesBook}) & Internal,
\Cref{S-thm:internal_combes_thomas}; conditional on the gap hypothesis (does not eliminate
HP gap input for the bounded-geometry case --- a standard locality tool for any future
internal bounded-geometry gap proof) \\
HP citation inventory and internalization roadmap (audit of every remaining HP citation in
both files, classified by type/status/difficulty; identifies the bounded-geometry closed-bulk
gap as the single research bottleneck) & Exhaustive line-by-line grep of both
\texttt{paper\_main.tex} and \texttt{paper\_supplement.tex}; no new mathematics & Audit-only,
\Cref{S-app:hp_inventory_iter12}; records the residual HP dependency precisely \\
Completeness verification for the seven internalized cases (per-theorem proof-chain
walk-through certifying that each of
\Cref{S-thm:internal_torus_gap,S-thm:internal_torus_chern,S-thm:internal_hex_gap,S-thm:internal_square_gap,S-thm:internal_hex_chern,S-thm:internal_square_chern,S-thm:internal_combes_thomas}
is genuinely Higson--Prodan-independent, modulo the OB-typed cap-product naming
convention for $S_K$) & Walk-through of each proof's transitive citation chain; classifies
each cited input as either (a) elementary closed-form computation in the manuscript or (b)
standard non-HP reference; closes the honesty gap that ``status I'' in
\Cref{S-tab:hp_inventory} could in principle have meant only that the line-grep inventory
did not see an HP citation & Audit-only, \Cref{S-app:completeness_iter13} and
\Cref{S-tab:completeness_iter13}; locks in the iter-1--6+11 HP-independence gains \\
Consolidated internal HP-replacement chapter (single navigation-friendly section gathering all
eight internal HP-replacement results --- the seven theorems
\Cref{S-thm:internal_torus_gap,S-thm:internal_torus_chern,S-thm:internal_hex_gap,S-thm:internal_square_gap,S-thm:internal_hex_chern,S-thm:internal_square_chern,S-thm:internal_combes_thomas}
plus the general operator-algebraic Chern convergence lemma
\Cref{S-thm:internal_bc_chern_general} --- under one unified-notation table, one dependency
roadmap, and one theorem-map table) & Cross-reference index pointing back to the per-iter
subsections that contain the full statements and proofs; reuses the iter-13 completeness
certifications; no new mathematics & Navigation-only,
\Cref{S-sec:internal_consolidated} and \Cref{S-tab:consolidated_internal_results}; reuses the
locked-in iter-13 HP-independence gains \\
Combes--Thomas verification (sub-step-by-sub-step audit of the iter-11 proof of
\Cref{S-thm:internal_combes_thomas}, verifying the four explicit constants
\(C_0(R,\Delta)\), \(\alpha_0(\gamma,\Lambda,R,\Delta)\), the contour resolvent bound
\(4/\gamma\), and \(C(\alpha,\gamma,\Lambda)=(\Lambda+\gamma/2)/\gamma\); confirms the
qualitative small-gap Combes--Thomas scaling
\(\alpha_0=O(\gamma/(\Lambda R))\); locates one Step~(d) sign-of-specialization
presentation slip that is silently repaired by the Hermitian symmetry of the self-adjoint
Fermi projection \(P_\pm(K)\)) & Sub-step-by-sub-step walk-through of the four
\((a)\)--\((d)\) sub-steps of the proof at supplement lines 7886--7996; per-sub-step
verification of the explicit constants and Combes--Thomas scaling; no new mathematics &
Audit-only, \Cref{S-app:combes_thomas_verification} and \Cref{S-tab:ct_audit_iter15};
independently verifies the iter-11 proof and constants \\
Pre-submission package audit (final integrity check of the manuscript: eight internal
HP-replacement results inventoried, every internal \(\backslash\)\texttt{Cref}, every
cross-document \texttt{xr} target, every \(\backslash\)\texttt{cite} key, every
\(\backslash\)\texttt{safeincludegraphics} file, and every \(\backslash\)\texttt{input}
table-include verified to resolve; one orphan citation \texttt{HigsonRoe} replaced by the
three-key form \texttt{HigsonRoeI,HigsonRoeII,HigsonRoeIII}; iter-15 Step~(d) caveat made
reachable from three independent vantage points) & Mechanical enumeration of labels,
citations, graphics, and input files; cross-document xr verification via the two
\texttt{.aux} files; no new mathematics & Audit-only,
\Cref{S-app:presubmission_audit_iter16},
\Cref{S-tab:presubmission_inventory_iter16},
\Cref{S-tab:presubmission_xref_audit_iter16},
\Cref{S-tab:presubmission_bib_audit_iter16},
\Cref{S-tab:presubmission_fig_audit_iter16}; certifies pre-submission integrity \\
Numerical-experiments framing audit (iter 18: per-figure/per-table audit of the
twenty-five numerical figures and tables in \Cref{S-sec:numerics} against the
iter-17 Part~A vs.\ Part~B carve-out
\Cref{S-tab:publishable_today}; verifies every figure/table is placed in Part~A
as a finite-matrix check of a Part~A theorem; verifies the four orientation
cross-references to Part~B theorems
[\(\protect\Cref{S-fig:edgeprop,S-fig:flux_spectral_flow,S-tab:flux_spectral_flow}\)
to \Cref{thm:pathb_roe_winding}; \(\protect\Cref{S-fig:refine}\) to
\Cref{thm:hp_imported}(i)] each already carry the canonical \emph{not-a-proof-input}
disclosure; certifies no figure or table is used as a proof input for any Part~A
theorem) & Mechanical per-caption read of every figure and table in
\Cref{S-sec:numerics}; cross-checked against the iter-17 framing
\Cref{S-app:publishable_today_iter17} and the iter-12 inventory
\Cref{S-app:hp_inventory_iter12}; no new mathematics & Audit-only,
\Cref{S-app:numerics_framing_audit_iter18},
\Cref{S-tab:numerics_framing_audit_iter18}; certifies that the iter-17 placement
of the numerical experiments in Part~A is consistent with the caption text \\
Supplement navigation index (iter 22: single navigation entry point at the start
of the supplement, listing every appendix subsection added in iters 1--21 plus
this iter-22 index itself, with subsection labels, one-line purposes, and page
references resolved automatically by \(\backslash\)\texttt{pageref}; eight rows
index the internal HP-replacement and topological-diagram content (iters 1--7,
10, 11), nine rows index the audit-, framing-, inventory-, consolidation-, and
review-content (iters 8, 9, 12--18, 21), and one row indexes the iter-22 index
itself; resolves the inherent tension between reader clarity and supplement bulk
by giving the reader a single entry point near the front of the document) &
Nineteen-row longtable indexing the appendix subsections added in iters 1--21;
no new mathematics, no recomputation, no figure regeneration & Audit-only,
\Cref{S-app:supplement_navigation_iter22},
\Cref{S-tab:supplement_navigation_iter22}; entry point for supplement
navigation \\
Quasi-uniform Delaunay first step (iter 23: define
\(\mathcal D(\theta_0,L)\), prove bounded-geometry parameter lemma,
state stability conjecture with proof outline) &
Angle-counting + Delaunay + cap-product; classical references
\cite{Cheeger1970,Chung1997,ChavelEigenvalues,Petersen2016,LottVillani2009,AdamsFournier,Strang1973}
for the four-ingredient roadmap & First step / partial /
conjectural at iter 23; conditional at iter 29 via iter-28 +
iter-29 cascade
(\Cref{S-rem:iter29_iter23_conditional_iter29});
\Cref{S-app:delaunay_first_step};
\Cref{S-lem:delaunay_bg_parameters} \emph{(proved)};
\Cref{S-thm:delaunay_gap_stability_conjecture}
\emph{(conjectural at iter 23; conditional at iter 29)};
\Cref{S-rem:delaunay_open_question}
\emph{(open)} \\
Cap-product coherence under quasi-isometric pullback (iter 24:
attempt at ingredient~(i) of iter-23's missing-ingredients list;
simplicial-isomorphism / translation sub-case proved rigorously;
vertex-perturbation extension proved; general QI case conjectural
with explicit support-and-norm decomposition) &
Functorial cap-product identity \cite[Section~3.3]{Hatcher};
simplicial chain map \cite[Ch.~1]{MunkresElements}; HP cap-product
construction \cite{HigsonProdanUniversalPRL,HigsonProdanUniversalArxiv};
Schur test \cite[Lem.~B.2.1]{ProdanSchulzBaldesBook} &
Sub-case proved / general case conjectural;
\Cref{S-app:cap_product_coherence_iter24};
\Cref{S-lem:translation_coherence_iter24} \emph{(proved)};
\Cref{S-lem:vertex_perturbation_coherence_iter24} \emph{(proved)};
\Cref{S-thm:qi_coherence_general_iter24}
\emph{(conjectural, proof outline)} \\
Hilbert-space isomorphism for different-dimension Delaunay pairs
(iter 25: attempt at ingredient~(ii) of iter-23's
missing-ingredients list; sub-complex inclusion sub-case proved
rigorously via inclusion partial isometry with explicit
boundary-supported correction; general common-refinement case
conjectural; combination with iter 24 would conditionally imply
iter-23 stability conjecture) &
Cap-product locality from HP construction
\cite{HigsonProdanUniversalPRL,HigsonProdanUniversalArxiv};
Schur test \cite[Lem.~B.2.1]{ProdanSchulzBaldesBook};
bounded-geometry constants from
\Cref{S-lem:delaunay_bg_parameters} &
Inclusion sub-case proved / general case conjectural at iter 25,
conditional at iter 29 via
\Cref{S-thm:three_step_composition_iter29};
\Cref{S-app:hilbert_iso_diff_dim_iter25};
\Cref{S-lem:inclusion_coherence_iter25} \emph{(proved)};
\Cref{S-thm:diff_dim_coherence_iter25}
\emph{(conjectural at iter 25; conditional at iter 29)} \\
Unified Delaunay-extension chapter overview (iter 26:
navigation/consolidation chapter integrating iters 23/24/25 into
a single chapter with status table, dependency diagram, reading
guide, and roadmap of open research items; no new mathematics,
no upgrade of conjectural results, no new bibliography entry) &
Source-order reorganization placing the iter-23/24/25
subsections inside a new overview section, plus a status table
reproducing the proof/conjecture/open status of each iter-23/24/25
result verbatim, a TikZ dependency diagram showing how the
iter-23 conjecture decomposes into ingredients~(i) and~(ii) and
how iters 24 and 25 each attack one of those ingredients, a
reading guide distinguishing math-content readers from
research-direction readers, and a six-item roadmap of open
research items (R1--R6) &
Navigation-only;
\Cref{S-sec:delaunay_extension_overview};
\Cref{S-tab:delaunay_status_iter26} \emph{(status table)};
\Cref{S-fig:delaunay_dependency_iter26}
\emph{(dependency diagram)} \\
Single-flip closed-form Schur norm bound (iter 27: closure of
sub-case~(b) of the iter-24 conjectural general QI theorem;
closed-form cap-product matrix elements on the eleven-simplex
local Hilbert space of a Delaunay flip quadrilateral; Schur
norm bound \(\|R_\phi^{\mathrm{single}}\|\le 4\) with the
explicit constant \(C_{\mathrm{flip}}=4\) independent of
\((\theta_0,L)\); sub-case~(c) multi-flip assembly remains
conjectural at iter 27) &
Closed-case cap-product convention
\cite[Section~3.3]{Hatcher}; HP cap-product locality
\cite{HigsonProdanUniversalPRL,HigsonProdanUniversalArxiv};
Schur test \cite[Lem.~B.2.1]{ProdanSchulzBaldesBook};
bounded-geometry constants from
\Cref{S-lem:delaunay_bg_parameters} &
Single-flip sub-case (b) proved internally; sub-case (c) still
conjectural at iter 27 (closed at iter 28);
\Cref{S-app:single_flip_closed_form_iter27};
\Cref{S-lem:cap_product_flip_iter27} \emph{(proved)};
\Cref{S-thm:single_flip_schur_iter27} \emph{(proved)} \\
Multi-flip Schur-test assembly (iter 28: closure of sub-case~(c)
of the iter-24 conjectural general QI theorem via telescoping
decomposition of the cumulative correction
\(R_\phi=V_\phi^*\,S_{K'}\,V_\phi-S_{K_0}\) and a
triangle-inequality bound
\(\|R_\phi\|\le C_{\mathrm{flip}}\,N_{\mathrm{flip}}=4N_{\mathrm{flip}}\);
edit-graph bounded-incidence lemma with explicit
\(D_*(\theta_0,L)=14\Delta_*(\theta_0)^2\); cascade upgrade of
\Cref{S-thm:qi_coherence_general_iter24} from conjectural to
proved internally with explicit constant
\(C_*(\theta_0,L)=C_{\mathrm{flip}}=4\) independent of
\((\theta_0,L)\), strictly tighter than the iter-24 provisional
\(8\,\Delta_*(\theta_0)\,L_*(\theta_0)\)) &
HP cap-product locality
\cite{HigsonProdanUniversalPRL,HigsonProdanUniversalArxiv};
Schur test \cite[Lem.~B.2.1]{ProdanSchulzBaldesBook};
bounded-geometry parameters from
\Cref{S-lem:delaunay_bg_parameters}; iter-27 per-flip
single-flip bound \Cref{S-thm:single_flip_schur_iter27}; iter-24
sub-case~(a) translation-coherence lemma
\Cref{S-lem:translation_coherence_iter24} &
Multi-flip sub-case (c) proved internally at iter 28; iter-24
conjectural theorem upgraded to proved via three-sub-case
cascade; \Cref{S-app:multi_flip_schur_iter28};
\Cref{S-lem:bounded_edit_incidence_iter28} \emph{(proved)};
\Cref{S-thm:multi_flip_schur_iter28} \emph{(proved)};
\Cref{S-thm:qi_coherence_general_iter24} \emph{(now proved at
iter 28 via cascade)} \\
\midrule
Common refinement for the Delaunay class (iter 29: candidate
construction for iter-25 sub-case~(c) generic common
refinement via the \emph{union-Delaunay refinement}
\(K_*=K_*(K,K')\), the Delaunay triangulation of the union
vertex set \(V_K\cup V_{K'}\); vertex inclusion
\(V_K\subseteq V_{K_*}\), \(V_{K'}\subseteq V_{K_*}\) proved;
parameter relaxation
\((\theta_0,L)\to(\theta_0',L')\) with explicit bounds
depending on the minimal inter-vertex distance
\(d_{\min}(V_K,V_{K'})\) and the maximum edge length
\(\overline{\ell}_{\max}\) proved; the union-Delaunay
refinement provides vertex inclusion but \emph{not} simplicial
sub-complex inclusion; conditional three-step composition
theorem closes iter-25 sub-case~(c) modulo a
sub-complex-refinement existence sub-conjecture, upgrading the
iter-23 stability conjecture from \textbf{conjectural} to
\textbf{conditional}) &
Closed-surface Delaunay convention from
\Cref{S-def:quasi_uniform_delaunay};
bounded-geometry parameters from
\Cref{S-lem:delaunay_bg_parameters}; iter-25 inclusion-coherence
lemma \Cref{S-lem:inclusion_coherence_iter25}; iter-28 cascade
\Cref{S-thm:qi_coherence_general_iter24}; HP cap-product
locality
\cite{HigsonProdanUniversalPRL,HigsonProdanUniversalArxiv};
Schur test \cite[Lem.~B.2.1]{ProdanSchulzBaldesBook} &
Union-Delaunay candidate construction for iter-25 sub-case~(c);
vertex-inclusion proved; sub-complex refinement is conditional;
iter-23 stability conjecture upgraded from conjectural to
conditional via iter-28 + iter-29 cascade;
\Cref{S-app:common_refinement_iter29};
\Cref{S-def:union_delaunay_iter29} \emph{(definition)};
\Cref{S-lem:union_delaunay_vertices_iter29} \emph{(proved)};
\Cref{S-lem:union_delaunay_parameters_iter29}
\emph{(proved)};
\Cref{S-rem:subcomplex_refinement_obstruction_iter29}
\emph{(honest, not proved)};
\Cref{S-thm:three_step_composition_iter29}
\emph{(conditional)};
\Cref{S-thm:delaunay_gap_stability_conjecture}
\emph{(conditional at iter 29 via the cascade)} \\
\midrule
Sub-complex refinement and weakened inclusion (iter 30:
negative-result iteration attacking the iter-29
sub-complex-refinement existence sub-conjecture via the
weakened \emph{refinement-in-the-union sense} relation; defines
a block partial isometry
\(V_K^\cup:\mathcal H(K)\hookrightarrow\mathcal H(\widetilde K_*)\)
for the constrained common refinement \(\widetilde K_*\) (the
overlay of \(K\) and \(K'\)), proves isometry, and explicitly
computes the cap-product matrix elements; identifies a concrete
renormalization mismatch factor \(\sqrt{N(\sigma')/N(\sigma)}\)
on incident pairs that produces a \emph{bulk-supported}
coherence correction with operator norm
\(\le 2\sqrt 2\Delta_*L_*\); the bulk support prevents direct
application of the iter-25 inclusion-coherence lemma in the
iter-29 three-step composition theorem; iter-23 stability
cascade status \textbf{unchanged from iter 29: conditional, not
unconditional}; documents the obstruction with an explicit
flat-torus example and outlines two future directions (weighted
Hilbert spaces and Roe-algebra coarse-geometry treatment) that
go beyond the scope of the iter-23--29 chain) &
Cap-product locality from HP construction
\cite{HigsonProdanUniversalPRL,HigsonProdanUniversalArxiv};
Schur test \cite[Lem.~B.2.1]{ProdanSchulzBaldesBook};
bounded-geometry parameters from
\Cref{S-lem:delaunay_bg_parameters}; iter-25 inclusion-coherence
lemma \Cref{S-lem:inclusion_coherence_iter25} (used as
counterpoint, not directly invoked); iter-29 union-Delaunay
construction \Cref{S-def:union_delaunay_iter29} (refined to the
overlay in \Cref{S-lem:overlay_refinement_iter30}) &
Negative-result iteration: iter-30 block-partial-isometry
construction is isometric but cap-product coherence correction
is bulk-supported, not boundary-supported; iter-23 stability
remains conditional at iter 29 (no upgrade);
\Cref{S-app:subcomplex_refinement_iter30};
\Cref{S-def:refinement_union_iter30} \emph{(definition)};
\Cref{S-lem:overlay_refinement_iter30} \emph{(proved)};
\Cref{S-def:block_partial_isometry_iter30} \emph{(definition)};
\Cref{S-lem:block_isometry_iter30} \emph{(proved)};
\Cref{S-lem:block_capproduct_matrix_iter30} \emph{(proved)};
\Cref{S-thm:block_obstruction_iter30} \emph{(proved
negative result)};
\Cref{S-rem:iter30_cascade_honest_iter30} \emph{(honest cascade
status: no upgrade from iter 29)} \\
\midrule
Cumulative Delaunay research chapter summary (iter 31:
navigation/consolidation chapter integrating iters 23--30 into
a single coherent presentation with three-tier classification
TIER 1 unconditional, TIER 2 conditional, TIER 3 negative
result; cumulative status table indexed by source iteration;
dependency-flow diagram showing how the foundational lemmas
of iters 23, 24, 25, 27, 28, 29, 30 feed into the iter-28
cascade closure (ingredient~(i) of iter-23) and the iter-29
conditional closure (ingredient~(ii) of iter-23), and how the
iter-30 negative result documents that approach~(A) does not
close the iter-29 sub-conjecture; roadmap of three remaining
open research directions (D1)--(D3); closing honest
assessment that the iter-23 stability conjecture closes
unconditionally for the periodic/Bravais flat-torus case via
iters 1--6 and is conditional in the general Delaunay class
on the single sub-complex refinement existence sub-conjecture
of iter 29; no new mathematics, no upgrade of any TIER 1, TIER
2, or TIER 3 status, no new bibliography entry) &
Source-order placement at the end of
\Cref{S-sec:delaunay_extension_overview}; cumulative status
table reproducing the TIER 1/2/3 classification of each
iter-23 to iter-30 result; TikZ dependency-flow diagram
showing the post-iter-30 dependency structure; cumulative
roadmap of three open research items integrating the (R1)--(R6)
items of iter 26 with the iter-28, iter-29, iter-30 closures
and obstructions; honest closing assessment documenting both
the unconditional flat-torus closure and the conditional
general-Delaunay closure &
Navigation-only;
\Cref{S-app:delaunay_chapter_summary_iter31};
\Cref{S-tab:delaunay_chapter_summary_iter31}
\emph{(cumulative status table)};
\Cref{S-fig:delaunay_chapter_summary_iter31}
\emph{(dependency-flow diagram)} \\
\midrule
Updated HP citation inventory (iter 32: re-grep both files for
every HP-citation occurrence after iters 13--31; classify each
cluster by type / status / iter-source; build a consolidated
representative-row inventory table; central honest finding that
\emph{no} HP closed-bulk case is newly internalized in iters
13--31; the bounded-geometry closed-bulk gap
\Cref{thm:hp_imported}(i) remains the single research bottleneck,
unchanged from iter 12; iters 13--22 audit / framing / navigation
cluster adds H-typed prose citations only; iters 23--31
Delaunay-extension cluster adds OB-typed cap-product locality
citations only, plus G-typed restatements of the iter-23 stability
conjecture target; no new mathematics, no upgrade of any TIER 1,
TIER 2, or TIER 3 status, no new bibliography entry) &
Exhaustive line-by-line grep of both \texttt{paper\_main.tex} and
\texttt{paper\_supplement.tex}; representative-row consolidation
of 485 raw HP-citation occurrences (90 in main + 395 in
supplement) into clusters indexed by source iteration; no new
mathematics &
Audit-only,
\Cref{S-app:hp_inventory_iter32},
\Cref{S-tab:hp_inventory_iter32}
\emph{(updated inventory table)};
\Cref{S-cor:no_new_hp_internalization_iters_13_31}
\emph{(central honest finding)};
\Cref{S-cor:single_bottleneck_iter32}
\emph{(bounded-geometry bottleneck reaffirmed)};
records the residual HP dependency at iter 32 \\
\midrule
Condensed internal development audit summary (iter 33:
paper-shortening pass relocating the audit/framing/navigation/
inventory cluster of iters 12, 13, 14, 16, 17, 18, 20, 22, 31, 32
into a single condensation table at the start of the supplement
with one row per audit iteration, columns iter / title /
one-sentence outcome / full subsection / approximate pages;
``Skip-on-first-reading notice'' paragraphs are added at the top of
each in-supplement audit/framing/navigation/inventory subsection;
``[skip-on-first-reading; iter-33 tag]'' tags are added to the
audit rows of the iter-22 supplement navigation index; effective
reader-relevant page count drops to \(\sim\!160\) pages of
load-bearing mathematics against a raw \(\sim\!246\)--\(248\) page
supplement; the iter-33 pass adds 1--3 pages and deletes nothing;
no new mathematics, no upgrade of any TIER 1, TIER 2, or TIER 3
status, no new bibliography entry, no modification of any internal
HP-replacement theorem) &
Source-order placement at the end of
\Cref{S-sec:reader_navigation}, immediately after the iter-22
supplement navigation index
\Cref{S-app:supplement_navigation_iter22}; ten-row condensation
table (iters 12, 13, 14, 16, 17, 18, 20, 22, 31, 32) with the
iter-20 row marked ``external'' because the iter-20 retrospective
is the external markdown file
\texttt{milestone\_retrospective\_iter20.md}; nine
``Skip-on-first-reading notice'' paragraph insertions; iter-22
navigation index ``one-line purpose'' column updates marking
audit/framing/navigation/inventory rows with the
``[skip-on-first-reading; iter-33 tag]'' tag; no new bibliography
entry &
Audit-only, paper-shortening pass,
\Cref{S-app:condensed_audit_summary_iter33},
\Cref{S-tab:condensed_audit_summary_iter33}
\emph{(condensation table)};
\Cref{S-lem:condensed_audit_summary_iter33_exposition}
\emph{(exposition-only certification)};
\Cref{S-cor:effective_page_count_iter33}
\emph{(post-iter-33 effective page count)};
entry point for audit-focused readers; resolves the supplement-%
bloat issue accumulated through iters 13--32 \\
\midrule
Final pre-submission audit (iter 35: post-iter-34 mechanical
integrity audit of the full submission package
\texttt{cover\_letter.tex} + \texttt{paper\_main.tex} +
\texttt{paper\_supplement.tex}; re-runs the iter-16
pre-submission audit
\Cref{S-app:presubmission_audit_iter16} against the
post-iter-34 manuscript state and extends it with three new
check categories absent at iter~16: the iter-19 stand-alone
cover letter, the iter-22 navigation index updated through
iter~33, and the iter-23--31 Delaunay-extension cluster; the
audit deliverable is a 21-row audit table with one row per
check item, columns audit item / status / notes; 20~rows
resolve to PASS, 1~row resolves to NOTE (the iter-22
navigation-index coverage gap for iters~23--34, a documented
editorial choice rather than a defect), 0~rows resolve to
NEEDS-FIX; pre-submission readiness statement records that the
post-iter-34 submission package is pre-submission ready
modulo the \texttt{[Journal name]} addressee-placeholder
replacement in the cover letter; no new mathematics, no
upgrade of any TIER 1, TIER 2, or TIER 3 status, no new
bibliography entry, no modification of any internal
HP-replacement theorem) &
Mechanical enumeration of labels, citations, graphics, and
\(\backslash\)\texttt{input} files; cross-document
\texttt{xr} verification via the two \texttt{.aux} files;
AI-prose-tell, Path-B/Path-C/scaffolding grep against
\texttt{cover\_letter.tex}, \texttt{paper\_main.tex},
\texttt{paper\_supplement.tex};
\(\backslash\)\texttt{section*\{Acknowledgments\}} and
\(\backslash\)\texttt{section*\{Declaration of generative AI
and AI-assisted technologies in the manuscript preparation
process\}} presence check; iter-33 skip-on-first-reading-notice
presence check; \texttt{check\_loops.py} validator pass; no new
mathematics &
Audit-only, final pre-submission integrity check,
\Cref{S-app:final_audit_iter35},
\Cref{S-tab:final_audit_iter35}
\emph{(21-row audit table)};
\Cref{S-lem:final_audit_iter35_exposition}
\emph{(exposition-only certification)};
\Cref{S-cor:final_audit_iter35_summary}
\emph{(twelve verifiable integrity conditions)};
verdict \textbf{pre-submission ready} \\
\midrule
Submission-package assembly (iter 36: clean
journal-submission zip assembled at
\texttt{doubling\_and\_jamming\_submission\_iter36.zip}
($\approx 7.02$~MB after DEFLATE compression) containing 51
files: 6 manuscript files
(\texttt{paper\_main.tex/pdf} + \texttt{paper\_supplement.tex/pdf}
+ \texttt{cover\_letter.tex/pdf}); 32 figure PNGs referenced by
\(\backslash\)\texttt{safeincludegraphics\{...\}} in the two
manuscripts; 9 table-input \texttt{.tex} files included via
\(\backslash\)\texttt{input\{...\}}; 3 companion-code files
(\texttt{periodic\_triangular\_bulk\_certificate.py} +
\texttt{run\_companion\_numerics\_with\_robustness.py} +
\texttt{requirements.txt}); and 1 \texttt{README.md} written
explicitly for the iter-36 package; explicit exclusions of the
555 \texttt{proof\_loop\_K} scaffolding directories, all
\texttt{.aux} / \texttt{.log} / \texttt{.out} / \texttt{.toc}
build artifacts, all \texttt{iterN\_summary.md} files, all
\texttt{DONE\_*.txt} markers, the helper scripts
\texttt{audit\_*.py} and \texttt{check\_loops.py}, the iter-1
split-source files \texttt{main.tex} and
\texttt{main\_split\_source.tex}, the iter-19 source-paper
exports, and the preserved snapshot directory; no new
mathematics, no theorem-status change, no new bibliography
entry, no modification of any internal HP-replacement theorem) &
Mechanical \texttt{Compress-Archive} packaging of the staged
deliverables into a flat zip; staging-directory file-count
verification; on-disk zip-size honest disclosure; recompile-from-
unzipped-directory acceptance criterion (zero
\texttt{LaTeX Warning}, zero undefined references, zero undefined
citations modulo the pre-existing typographic warnings); no new
mathematics &
Audit-only, mechanical packaging step,
\Cref{S-app:submission_package_iter36}
\emph{(submission-package assembly subsection)};
\Cref{S-lem:submission_package_iter36_inventory}
\emph{(exposition-only inventory)};
\Cref{S-thm:submission_package_iter36_closure}
\emph{(packaging-closure theorem)};
entry point for the journal-submission step; verdict
\textbf{clean submission zip ready}, modulo the
\texttt{[Journal name]} addressee-placeholder replacement at
line~36 of \texttt{cover\_letter.tex} \\
\bottomrule
\end{tabular}
\end{table}

\begin{figure}[H]
\centering
\[
\begin{tikzcd}[row sep=2.6em, column sep=1.4em]
|[draw, rounded corners, fill=gray!18]|
  \begin{array}{c}\text{\textbf{HP gap input}}\\[-0.2em]\text{\Cref{thm:hp_imported}(i)}\end{array}
  \arrow[r, dashed, no head]
& |[draw, rounded corners, fill=gray!18]|
  \begin{array}{c}\text{\textbf{HP Chern input}}\\[-0.2em]\text{\Cref{thm:hp_imported}(ii)}\end{array}
  \arrow[r, "\text{closed-bulk}"]
& |[draw, rounded corners, fill=white]|
  \begin{array}{c}\text{boundary map}\\[-0.2em]\text{\(\partial_Y\): bulk \(\to\) interface}\end{array}
  \arrow[r, "\text{Toeplitz}"]
& |[draw, rounded corners, fill=white]|
  \begin{array}{c}\text{interface \(K_1\)-class}\\[-0.2em]\text{winding \(\pm1\)}\end{array} \\
|[draw, rounded corners, fill=green!22, double, double distance=2pt]|
  \begin{array}{c}\text{\textbf{Flat-torus gap}}\\[-0.2em]\text{\Cref{S-thm:internal_torus_gap}}\end{array}
  \arrow[u, "\text{replaces (gap)}"', dashed]
& |[draw, rounded corners, fill=green!22, double, double distance=2pt]|
  \begin{array}{c}\text{\textbf{Flat-torus Chern}}\\[-0.2em]\text{\Cref{S-thm:internal_torus_chern}}\end{array}
  \arrow[u, "\text{replaces (Chern, op.\ level)}"', dashed]
& |[draw, rounded corners, fill=green!12]|
  \begin{array}{c}\text{Periodic Toeplitz}\\[-0.2em]\text{\Cref{thm:periodic_toeplitz_pairing}}\end{array}
  \arrow[u, "\text{(internal)}"', dashed]
& |[draw, rounded corners, fill=green!12]|
  \begin{array}{c}\text{Internal winding\,\(\pm1\)}\\[-0.2em]\text{\Cref{cor:toeplitz_winding}}\end{array}
  \arrow[u, "\text{(internal)}"', dashed]
\\
|[draw, rounded corners, fill=green!22, double, double distance=2pt]|
  \begin{array}{c}\text{\textbf{Hexagonal gap}}\\[-0.2em]\text{\Cref{S-thm:internal_hex_gap}}\end{array}
  \arrow[uu, "\text{replaces (gap, hex)}"', dashed, bend right=30]
& |[draw, rounded corners, fill=green!22, double, double distance=2pt]|
  \begin{array}{c}\text{\textbf{Hexagonal Chern}}\\[-0.2em]\text{\Cref{S-thm:internal_hex_chern}}\end{array}
  \arrow[uu, "\text{replaces (Chern, hex)}"', dashed, bend right=15]
& |[draw, rounded corners, fill=green!12]|
  \begin{array}{c}\text{Periodic triangular Chern (symbol)}\\[-0.2em]\text{\Cref{thm:periodic_triangular_chern}}\end{array}
  \arrow[ul, "\text{FHS limit}"', dashed]
& \\
|[draw, rounded corners, fill=green!22, double, double distance=2pt]|
  \begin{array}{c}\text{\textbf{Square gap}}\\[-0.2em]\text{\Cref{S-thm:internal_square_gap}}\end{array}
  \arrow[uuu, "\text{replaces (gap, sq)}"', dashed, bend left=15]
& |[draw, rounded corners, fill=green!22, double, double distance=2pt]|
  \begin{array}{c}\text{\textbf{Square Chern}}\\[-0.2em]\text{\Cref{S-thm:internal_square_chern}}\end{array}
  \arrow[uuu, "\text{replaces (Chern, sq)}"', dashed, bend left=15]
& & \\
&
|[draw, rounded corners, fill=green!12]|
  \begin{array}{c}\text{\textbf{General lattice-Chern lemma}}\\[-0.2em]
  \text{\Cref{S-thm:internal_bc_chern_general}}\end{array}
  \arrow[u, "\text{specializes to (sq)}"', dashed]
  \arrow[uu, "\text{specializes to (hex)}", dashed, bend left=18]
  \arrow[uuu, "\text{specializes to (tri)}", dashed, bend left=30]
& & \\
|[draw, rounded corners, fill=green!12]|
  \begin{array}{c}\text{\textbf{Internal Combes--Thomas}}\\[-0.2em]
  \text{\Cref{S-thm:internal_combes_thomas}}\end{array}
  \arrow[r, dashed, "\text{tool}"]
& |[draw, rounded corners, fill=gray!8, dashed]|
  \begin{array}{c}\text{Bounded-geometry case}\\[-0.2em]\text{(future iter)}\end{array}
  \arrow[uuuuu, dashed, bend right=30, "\text{future BG}"']
& & \\
\end{tikzcd}
\]
\par\medskip
\noindent\fbox{\begin{minipage}{0.96\linewidth}\centering\footnotesize
\textbf{Legend.}\
\tikz[baseline=-0.5ex]\node[draw, rounded corners, fill=green!22, double, double distance=1.5pt,
inner sep=2pt]{};\ doubled green border: internal HP-replacement (this paper).\quad
\tikz[baseline=-0.5ex]\node[draw, rounded corners, fill=gray!18, inner sep=2pt]{};\ grey:
external HP citation (still required for general bounded-geometry case).\quad
\tikz[baseline=-0.5ex]\node[draw, rounded corners, fill=white, inner sep=2pt]{};\ white:
downstream construction (boundary map, interface class, winding).\quad
\tikz[baseline=-0.5ex]\node[draw, rounded corners, fill=green!12, inner sep=2pt]{};\ light
green (single border): internal symbol-level input feeding an internal replacement above.
\textbf{Columns:} (1) gap inputs; (2) Chern inputs; (3)--(4) downstream construction (boundary
map, interface \(K_1\)-class, integer winding). \textbf{Rows:} (top) HP and downstream
construction; (row 2) triangular flat-torus internal replacements; (row 3) hexagonal benchmark
internal replacements; (row 4) square benchmark internal replacements; (row 5) the single
internal-replacement node ``General lattice-Chern lemma'' (\Cref{S-thm:internal_bc_chern_general})
that subsumes the three Chern-column nodes above as corollaries; (row 6) the internal
Combes--Thomas exponential decay tool (\Cref{S-thm:internal_combes_thomas}) feeding a
dashed-grey placeholder ``Bounded-geometry case (future iter)'' that records the
remaining open HP-replacement target. The single dashed-grey placeholder is the only non-internal
node in rows 2--6.
\end{minipage}}
\caption{Dependency structure for the closed-bulk-to-interface chain, restructured into a
clean four-column layout. \textbf{Columns:} column 1 collects \emph{gap} inputs; column 2
collects \emph{Chern} inputs; columns 3--4 collect the downstream construction (boundary map
\(\partial_Y\), interface \(K_1\)-class, integer winding). \textbf{Rows:} the top row contains
the HP inputs (grey, external) on the left and the downstream construction (white) on the
right. The three lower rows are internal HP-replacements organized by benchmark family: row 2
holds the triangular flat-torus replacements (gap and Chern, on the universal Hamiltonian),
row 3 holds the hexagonal benchmark replacements (gap and rank-\(1\) Chern), and row 4 holds
the square benchmark replacements (gap and rank-\(1\) Chern). The single-bordered green node
``Periodic triangular Chern (symbol)'' (\Cref{thm:periodic_triangular_chern}) in row 3,
column 3 is the symbol-level Chern input that feeds the operator-level row-2 ``Flat-torus
Chern'' node via the FHS limit. The single-bordered green node ``General lattice-Chern lemma''
(\Cref{S-thm:internal_bc_chern_general}) in row 5, column 2 is the operator-algebraic Chern
convergence statement for arbitrary continuously parametrized self-adjoint Bloch families on
\(\mathbb T^2\); the three upward dashed arrows from this node mark the three corollaries
(triangular, hexagonal, square) that recover the row-2/3/4 Chern nodes by specializing to the
benchmark-specific data. The single-bordered green
\Cref{thm:periodic_toeplitz_pairing} and \Cref{cor:toeplitz_winding} were internal already.
After this iteration, the gap and Chern inputs of HP are fully internal on three benchmarks
(triangular flat torus, hexagonal benchmark on the same flat torus, and square benchmark on
the same flat torus); only the general bounded-geometry case retains a dependency on the
top-row HP nodes. The new row-6 ``Internal Combes--Thomas''
(\Cref{S-thm:internal_combes_thomas}) is a standard locality tool: given a uniform Fermi gap
on a bounded-geometry triangulation, it gives uniform off-diagonal exponential decay of the
Fermi projection with a rate independent of the underlying triangulation; the dashed
``will replace (BG)'' arrow indicates that this decay is a prerequisite for any future
internal proof of the bounded-geometry case but does not by itself supply the bounded-geometry
gap. The legend box below the diagram defines the colour and border
conventions. The internal compatibility of the Roe connecting morphism \(\partial_Y\) of
column 3 with the Ewert--Meyer Mayer--Vietoris boundary map is the content of
\Cref{thm:pathb_roe_winding}, visualized as a separate SES-comparison diagram
(\Cref{fig:roe_em_ses_comparison}) immediately before its body proof.}
\label{fig:hp_dependency_diagram}
\end{figure}

\begin{table}[H]
\centering
\small
\caption{Notation used for the bulk, jammed, and sidewise interface operators.}
\label{tab:notation_ledger}
\begin{tabular}{>{\raggedright\arraybackslash}p{0.24\linewidth}
                >{\raggedright\arraybackslash}p{0.65\linewidth}}
\toprule
Symbol & Meaning \\
\midrule
\(H_\pm=(B+B^\dagger)\pm S\) & Closed universal Chern Hamiltonian. \\
\(H_{\pm,M}=H_\pm(D(\K))+M\Proj_{\mathrm{copy}}\) & Finite doubled-and-jammed Hamiltonian. \\
\(A=P_{\mathrm{orig}}H_\pm(D(\K))P_{\mathrm{orig}}\) & Original-side compression of the doubled
Hamiltonian. \\
\(H_\pm^{\mathrm{PL}}(\K)\) & Intrinsic Poincar\'e--Lefschetz boundary Hamiltonian. \\
\(P_\pm=\chi_{(-\infty,0)}(H_\pm^\infty)\) & Closed-bulk Fermi projection in the coarse disjoint
union. \\
\(J_{\pm,M}^\infty\) & Direct-sum jammed family defining the interface. \\
\(T^+\), \(T^-\) & Positive-side and negative-side comparison operators. \\
\(Y\), \(C^*(Y\subset X)\) & Coarse interface and its Roe ideal. \\
\(\partial_Y(E_+(\pi(P_\pm)))\) & Internally constructed interface \(K_1\)-class. \\
\bottomrule
\end{tabular}
\end{table}

%====================================================================
\section{Simplicial setup and universal operators}\label{sec:setup}
%====================================================================

\begin{definition}[Triangulated surface with boundary]\label{def:surfwithbdry}
A finite oriented $2$-dimensional simplicial complex $\K$ is a \emph{triangulated surface with
boundary} if its geometric realization $|\K|$ is a compact PL $2$-manifold with boundary and the
boundary subcomplex $\bd\K$ is a \emph{full subcomplex}: every simplex of $\K$ whose vertices all
lie on $\bd\K$ is itself a boundary simplex. Every edge is incident to exactly one or two faces;
the edges incident to exactly one face form $\bd\K$, and $\bd\K$ is a disjoint union of cycles.
\end{definition}

The fullness condition is mild: it holds automatically after one barycentric subdivision. Its
role in keeping the simplicial double an abstract simplicial complex (without duplicate simplices
or boundary-edge valence failures) is discussed in detail in the supplement,
Section~\ref{S-sec:setup_remarks}.

We work on the simplicial Hilbert space
\(\Hil(\K)=\bigoplus_{p=0}^2 \ell^2(\K_p)\)
with the canonical orthonormal basis indexed by oriented simplices. The simplicial boundary
operator $B$ is the standard chain boundary, viewed as a degree-lowering operator. Then
$B^2=0$ and $B+B^\dagger$ is self-adjoint and local. We endow the set of barycenters
$X(\K)$ with the graph metric on the $1$-skeleton of the first barycentric subdivision; an
operator $T$ has \emph{propagation at most $R$} when
$\langle\sigma,T\tau\rangle=0$ for $\dist(\sigma,\tau)>R$.

For a \emph{closed} oriented triangulation $\K$, cap product with the fundamental cycle
$|X\rangle=\sum_{\tau\in\K_2}|\tau\rangle$ gives the Higson--Prodan operator $P$ from which one
defines a Hermitian, degree-reversing, local operator $S$ by the degree-dependent
symmetrization
\[
T|_{\Hil_p}=\tfrac12\bigl(P^\dagger+(-1)^pP\bigr),\qquad
S|_{\Hil_p}=i^{p(p-1)+1}T|_{\Hil_p},
\]
so that $BS+SB^\dagger=0$. The closed universal Chern Hamiltonian is
\(H_\pm(\K)=(B+B^\dagger)\pm S\).

\subsection{Closed-case gap criteria}
Two finite-dimensional gap criteria are used as building blocks; their proofs are short and
collected in the supplement.

\begin{proposition}[Exact gap lower bound under exact anticommutation]\label{prop:gap_from_S}
Let $\K$ be closed and assume $(B+B^\dagger)S+S(B+B^\dagger)=0$. Then
$H_\pm(\K)^2=(B+B^\dagger)^2+S^2$. If $\mu_S:=\dist(0,\spec(S))>0$, then
$\spec(H_\pm(\K))\cap(-\mu_S,\mu_S)=\emptyset$.
\end{proposition}
\proofsuppref{app:proof:prop:gap_from_S}

\begin{proposition}[Gap lower bound with anticommutation defect]\label{prop:gap_defect}
With $Q=B+B^\dagger$, $\mu_S:=\dist(0,\spec(S))$, and $\e_{\mathrm{ac}}:=\opnorm{QS+SQ}$, if
$S$ is Hermitian then $H_\pm(\K)^2\succeq(\mu_S^2-\e_{\mathrm{ac}})\Id$.
\end{proposition}
\proofsuppref{app:proof:prop:gap_defect}

\subsection{Doubling, jamming, and the boundary Hamiltonians}
Given $\K$ with boundary in the sense of Definition~\ref{def:surfwithbdry}, let $\overline\K$ be an
oriented copy with reversed face orientation, sharing each boundary vertex of $\K$. The
\emph{double} $D(\K)=\K\cup_{\bd\K}\overline\K$ is a finite closed oriented triangulated surface:

\begin{proposition}[The double is a closed oriented triangulated surface]\label{prop:double_closed}
Let $\K$ satisfy Definition~\ref{def:surfwithbdry}. Then $D(\K)$ is a finite closed oriented
abstract simplicial complex that triangulates a closed surface.
\end{proposition}
\proofsuppref{app:proof:prop:double_closed}

\begin{theorem}[Simplex count identities under doubling]\label{thm:counts}
\(|D(\K)_2|=2|\K_2|\),\, \(|D(\K)_1|=2|\K_1|-|\bd\K_1|\),\, \(|D(\K)_0|=2|\K_0|-|\bd\K_0|\), and
$|\bd\K_0|=|\bd\K_1|$. Consequently
$\dim\Hil(D(\K))=2\dim\Hil(\K)-2|\bd\K_0|$, $\chi(D(\K))=2\chi(\K)$, and for a connected,
orientable, genus-$g$ surface with $b\ge 1$ boundary components, the doubled genus is
$g_{\mathrm{double}}=2g+b-1$.
\end{theorem}
\proofsuppref{app:proof:thm:counts}

The closed-case construction therefore produces the operator $H_\pm(D(\K))$ on $\Hil(D(\K))$.
We split
\[
\Hil(D(\K))=\Hil_{\mathrm{orig}}\oplus\Hil_{\mathrm{copy}}
\]
and \emph{jam} the reflected copy by the on-site potential
\begin{equation}\label{eq:Hjam}
H_{\pm,M}=H_\pm(D(\K))+M\Proj_{\mathrm{copy}}.
\end{equation}
The block decomposition
\(H_\pm(D(\K))=\begin{pmatrix}A&C\\ C^\dagger&D\end{pmatrix}\)
on $\Hil_{\mathrm{orig}}\oplus\Hil_{\mathrm{copy}}$ defines the
\emph{original-side boundary Hamiltonian}
\(H_{\pm}^{\bd}(\K):=A\), a self-adjoint local operator on the surface-with-boundary.

The intrinsic alternative uses the relative chain complex
$C_p(\K,\bd\K)=C_p(\K)/C_p(\bd\K)$ and the relative fundamental chain
$[\K,\bd\K]=\sum_{\tau\in\K_2}|\tau\rangle$. Cap product with $[\K,\bd\K]$, followed by the
canonical lift $\widehat P_p:=j_{2-p}P^{\mathrm{rel}}_p:C_p(\K)\to C_{2-p}(\K)$, gives the
Hermitian, degree-reversing, local relative duality operator
$S^{\mathrm{PL}}_\K$ on $\Hil(\K)$. The \emph{relative Poincar\'e--Lefschetz boundary
Hamiltonian} is
\begin{equation}\label{eq:H_PL}
H^{\mathrm{PL}}_{\pm}(\K):=B_\K+B_\K^\dagger\pm S^{\mathrm{PL}}_\K\quad\text{on }\Hil(\K).
\end{equation}

\begin{lemma}[Relative cap identity and collar defect]\label{lem:relative_cap_identity}
With $\partial^{\mathrm{rel}}_r:C_r(\K,\bd\K)\to C_{r-1}(\K,\bd\K)$ defined by
$\partial^{\mathrm{rel}}_r q_r=q_{r-1}B_{\K,r}$, and $\delta_p:=B_{\K,p+1}^\dagger$, the
relative cap-product blocks satisfy
$\partial^{\mathrm{rel}}_{2-p}P^{\mathrm{rel}}_p+(-1)^p P^{\mathrm{rel}}_{p+1}\delta_p=0$
for $p=0,1$. After lifting by $j$, the difference
\(B_{\K,2-p}\widehat P_p+(-1)^p\widehat P_{p+1}\delta_p\) has matrix coefficients only when
both source and target lie in a fixed barycentric collar of $\bd\K$. Consequently the relative
duality defect $B_\K S^{\mathrm{PL}}_\K+S^{\mathrm{PL}}_\K B_\K^\dagger$ is supported in a fixed
boundary collar.
\end{lemma}
\proofsuppref{app:proof:lem:relative_cap_identity}

\begin{lemma}[Coefficient agreement off the boundary collar]\label{lem:cap_coefficient_agreement}
Let $P_D$ be the closed cap-product operator on $D(\K)$ and $\widehat P$ the lifted relative
cap-product operator on $\K$. There is a radius $R_\frown$ depending only on the local
cap-product rule, such that
$\langle\sigma,P_D\tau\rangle=\langle\sigma,\widehat P\tau\rangle$ whenever
$\sigma,\tau\in\K\setminus N_{R_\frown}(\bd\K)$.
\end{lemma}
\proofsuppref{app:proof:lem:cap_coefficient_agreement}

\begin{theorem}[Relative boundary Hamiltonian and doubled compression]\label{thm:relative_boundary_comparison}
Let \(\K\) be a triangulated oriented surface with full boundary subcomplex. Under the natural
identification of \(\Hil(\K)\) with \(\Hil_{\mathrm{orig}}\subset\Hil(D(\K))\),
\(H^{\mathrm{PL}}_{\pm}(\K)\) is self-adjoint and local. Moreover there is a self-adjoint operator
\(R_{\partial,\pm}\) such that
\begin{equation}\label{eq:relative_boundary_comparison}
H_{\pm}^{\bd}(\K)
=
H^{\mathrm{PL}}_{\pm}(\K)+R_{\partial,\pm},
\end{equation}
and the comparison error \(R_{\partial,\pm}\), and also the relative duality defect
\(B_\K S^{\mathrm{PL}}_\K+S^{\mathrm{PL}}_\K B_\K^\dagger\), are supported in a bounded collar of
the boundary: there is a radius \(R_{\mathrm{PL}}\), depending only on the local cap-product
convention, such that the matrix entries vanish unless both source and target lie in
\(N_{R_{\mathrm{PL}}}(\bd\K)\). For uniformly regular refinement families the same
\(R_{\mathrm{PL}}\) works for all members.
\end{theorem}

\paragraph{Proof outline of \Cref{thm:relative_boundary_comparison}.}
The proof has four steps; the full version is in
\Cref{S-app:proof:thm:relative_boundary_comparison}.
(i) Self-adjointness and locality of \(H^{\mathrm{PL}}_\pm(\K)\) follow from the construction of
the relative duality operator \(S^{\mathrm{PL}}_\K\) (cap product with \([\K,\bd\K]\) followed
by the canonical lift \(j_{2-p}P^{\mathrm{rel}}_p\)); the cap-product rule is local and the lift
\(j\) is the standard quotient inclusion.
(ii) The matrix-element comparison of \(P_D\) (closed double cap product) and
\(\widehat P\) (lifted relative cap product) is the content of
\Cref{lem:cap_coefficient_agreement}: away from a fixed barycentric collar of \(\bd\K\) the
two cap-product blocks have identical matrix entries by the local nature of the simplicial
cap-product formula.
(iii) Consequently, the difference
\(R_{\partial,\pm}:=H_\pm^{\bd}(\K)-H^{\mathrm{PL}}_\pm(\K)\) has matrix entries supported in
the collar \(N_{R_{\mathrm{PL}}}(\bd\K)\), where \(R_{\mathrm{PL}}\) is the maximum
propagation of the local cap-product rule plus the maximum propagation of the lift.
(iv) The relative duality defect
\(B_\K S^{\mathrm{PL}}_\K+S^{\mathrm{PL}}_\K B_\K^\dagger\) is supported in the same collar by
\Cref{lem:relative_cap_identity}: the relative cap-product blocks satisfy
\(\partial^{\mathrm{rel}}_{2-p}P^{\mathrm{rel}}_p+(-1)^p P^{\mathrm{rel}}_{p+1}\delta_p=0\) on
the quotient \(C_\bullet(\K,\bd\K)\), so the lifted version has nonzero contributions only at
the relative boundary, supported in the same fixed collar. For uniformly regular refinement
families the local cap-product rule is the same, so the same \(R_{\mathrm{PL}}\) works for all
members.

\proofsuppref{app:proof:thm:relative_boundary_comparison}

\begin{lemma}[Spectral counts under collar-supported perturbations]\label{lem:collar_rank_counts}
If $R=P_W R P_W$ for a projection $P_W$ of rank $r$, then
$|N_{A+R}(t)-N_A(t)|\le r$ for all $t$ and
$|\rank\chi_I(A+R)-\rank\chi_I(A)|\le 2r$ for every interval $I$. Applied to
$A=H^{\mathrm{PL}}_{\pm}(\K)$ and $A+R=H_{\pm}^{\bd}(\K)$, one may take
$W=\ell^2(N_{R_{\mathrm{PL}}}(\bd\K))$. In particular, along refinement families with
$r_n/d_n\to0$, the normalized spectral distribution functions of the two boundary models share
the same subsequential limits.
\end{lemma}
\proofsuppref{app:proof:lem:collar_rank_counts}

%====================================================================
\section{Finite-dimensional decoupling theory}\label{sec:decoupling}
%====================================================================

The finite-dimensional analysis is perturbation theory specialized to the block structure produced
by doubling and jamming. We use standard Schur-complement and spectral projection tools
\cite[Chapter~II]{KatoPerturbation,BhatiaMatrixAnalysis}, but keep the constants and separation
hypotheses explicit because they are later checked on the computed matrices. The full proofs are
collected in the online supplement, Section~\ref{S-sec:schur_proofs}.

\begin{lemma}[Schur complement resolvent formula]\label{lem:schur}
Let $T=\bigl(\begin{smallmatrix}A&C\\ C^\dagger&D\end{smallmatrix}\bigr)$ be self-adjoint on
$\Hil_1\oplus\Hil_2$ and $z\in\mathbb C\setminus\mathbb R$. Then the $(1,1)$ block of $(T-z)^{-1}$
is $\bigl(A-z-C(D-z)^{-1}C^\dagger\bigr)^{-1}$.
\end{lemma}

Write $D_M:=D+M\Id$ and $d_M(z):=\dist(z,\spec(D_M))$.

\begin{theorem}[Compressed resolvent convergence with explicit $O(1/M)$ rate]\label{thm:resolvent}
Let $H_{\pm,M}$ be as in \eqref{eq:Hjam}. For any $z\in\mathbb C\setminus\mathbb R$,
\begin{equation}\label{eq:resolvent_bound_exact}
\opnorm{P_{\mathrm{orig}}(H_{\pm,M}-z)^{-1}P_{\mathrm{orig}}-(A-z)^{-1}}
\le
\frac{\opnorm{C}^2}{d_M(z)\,|\Im z|^2}.
\end{equation}
In particular, if $M>\opnorm{D}+|\Re z|$,
$P_{\mathrm{orig}}(H_{\pm,M}-z)^{-1}P_{\mathrm{orig}}\to(A-z)^{-1}$ in operator norm at rate
$O(1/M)$, so the doubled-and-jammed Hamiltonian recovers the original-side boundary Hamiltonian
$A=H_{\pm}^{\bd}(\K)$ in the norm-resolvent sense as $M\to\infty$.
\end{theorem}

\begin{proof}[Proof sketch; full proof in Supplement, Section~\ref{S-sec:schur_proofs}.]
Write $R_M:=P_{\mathrm{orig}}(H_{\pm,M}-z)^{-1}P_{\mathrm{orig}}$ and $R_0:=(A-z)^{-1}$.
The Schur complement (Lemma~\ref{lem:schur}) gives the resolvent identity
$R_M-R_0=R_M\,K_M\,R_0$ with the self-energy correction
$K_M:=C(D_M-z)^{-1}C^\dagger$. The resolvent norms satisfy
$\opnorm{R_M}\le|\Im z|^{-1}$ and $\opnorm{R_0}\le|\Im z|^{-1}$, and the self-energy is
bounded by $\opnorm{K_M}\le\opnorm{C}^2/d_M(z)$. Multiplying gives
\eqref{eq:resolvent_bound_exact}.
\end{proof}

\begin{corollary}[Boundary Hamiltonian as the large-jam limit]\label{cor:boundary_hamiltonian_limit}
Let $H_{\pm}^{\bd}(\K)=A$ be the original-side boundary Hamiltonian. For every
$z\in\mathbb C\setminus\mathbb R$,
\(P_{\mathrm{orig}}(H_{\pm,M}-z)^{-1}P_{\mathrm{orig}}\to(H_{\pm}^{\bd}(\K)-z)^{-1}\) in operator
norm with the explicit $O(M^{-1})$ bound \eqref{eq:resolvent_bound_exact}. Thus the
doubled-and-jammed Hamiltonian recovers a self-adjoint boundary Hamiltonian on the original side
in the norm-resolvent sense; equivalently, by Theorem~\ref{thm:relative_boundary_comparison},
the same limit is $(H^{\mathrm{PL}}_{\pm}(\K)+R_{\partial,\pm}-z)^{-1}$.
\end{corollary}
\proofsuppref{app:proof:cor:boundary_hamiltonian_limit}

\begin{corollary}[Degenerate full block-resolvent limit]\label{cor:full_resolvent}
For $z\in\mathbb C\setminus\mathbb R$, $(H_{\pm,M}-z)^{-1}\to(A-z)^{-1}\oplus0$ in operator norm
at rate $O(1/M)$, with the explicit triangle inequality
\(\|(H_{\pm,M}-z)^{-1}-((A-z)^{-1}\oplus0)\|\le \opnorm{C}^2/(d_M(z)|\Im z|^2)
+ 2\opnorm{C}/(d_M(z)|\Im z|)+1/d_M(z)+\opnorm{C}^2/(d_M(z)^2|\Im z|)\).
This is generalized norm-resolvent convergence in which the copy sector is sent to infinite energy.
\end{corollary}
\proofsuppref{app:proof:cor:full_resolvent}

\begin{theorem}[First-order Schur-complement expansion]\label{thm:first_order_eff}
Fix $z\in\mathbb C\setminus\mathbb R$. If $M>\opnorm{D-z}$ and $A_M^{(1)}:=A-\frac1MCC^\dagger$,
then
\begin{equation}\label{eq:first_order_bound}
\opnorm{P_{\mathrm{orig}}(H_{\pm,M}-z)^{-1}P_{\mathrm{orig}}-(A_M^{(1)}-z)^{-1}}
\le
\frac{\opnorm{C}^2\,\opnorm{D-z}}{M(M-\opnorm{D-z})\,|\Im z|^2}=O(M^{-2}).
\end{equation}
\end{theorem}
\proofsuppref{app:proof:thm:first_order_eff}

For a positively oriented simple closed contour $\Gamma\subset\rho(T)$ write the Riesz projection
$P_\Gamma(T):=\frac{1}{2\pi i}\int_\Gamma (z-T)^{-1}\,dz$. Set
$\delta_A(\Gamma):=\dist(\Gamma,\spec(A))$,
$\delta_D(\Gamma):=\dist(\Gamma,\spec(D_M))$, and
$\delta_H(\Gamma):=\delta_A(\Gamma)-\opnorm{C}^2/\delta_D(\Gamma)$.

\begin{theorem}[General contour-projection comparison]\label{thm:proj_general}
Assume $\delta_A(\Gamma)>0$, $\delta_D(\Gamma)>0$, $\delta_H(\Gamma)>0$. Then
$\Gamma\subset\rho(H_{\pm,M})$,
$\rank P_\Gamma(H_{\pm,M})=\rank P_\Gamma(A)+\rank P_\Gamma(D_M)$, and
\begin{equation}\label{eq:proj_general_bound}
\opnorm{P_\Gamma(H_{\pm,M})-\bigl(P_\Gamma(A)\oplus P_\Gamma(D_M)\bigr)}
\le \tfrac{|\Gamma|}{2\pi}\,\eta_M(\Gamma)=O(M^{-1}),
\end{equation}
where $\eta_M(\Gamma)$ is the explicit constant
\[
\eta_M(\Gamma)
=\frac{\opnorm C^2}{\delta_A(\Gamma)\delta_H(\Gamma)\delta_D(\Gamma)}
+\frac{2\opnorm C}{\delta_H(\Gamma)\delta_D(\Gamma)}
+\frac{\opnorm C^2}{\delta_H(\Gamma)\delta_D(\Gamma)^2}.
\]
\end{theorem}
\proofsuppref{app:proof:thm:proj_general}

For $E_{\max}>0$ define $g_M(E_{\max}):=\dist\bigl([-E_{\max},E_{\max}],\spec(D_M)\bigr)$.

\begin{theorem}[Bounded-energy spectral localization]\label{thm:spectral_localization}
If $E$ is an eigenvalue of $H_{\pm,M}$ with $|E|\le E_{\max}$ and $g_M(E_{\max})>0$, then
$\dist(E,\spec(A))\le\opnorm{C}^2/g_M(E_{\max})$.
\end{theorem}
\proofsuppref{app:proof:thm:spectral_localization}

\begin{corollary}[Spectral comparison with the relative boundary Hamiltonian]\label{cor:relative_spectral_localization}
Under the hypotheses of Theorem~\ref{thm:spectral_localization},
\(\dist\bigl(E,\spec(H^{\mathrm{PL}}_\pm(\K)+R_{\partial,\pm})\bigr)
\le\opnorm{C}^2/g_M(E_{\max})\).
\end{corollary}
\proofsuppref{app:proof:cor:relative_spectral_localization}

\begin{theorem}[Bounded-energy eigenvalues lie near the first-order effective operator]\label{thm:first_order_spectral_localization}
Let $E$ be an eigenvalue of $H_{\pm,M}$ with $|E|\le E_{\max}$ and $g_M(E_{\max})>0$. Then
\begin{equation}\label{eq:specA1_exact}
\dist\bigl(E,\spec(A_M^{(1)})\bigr)
\le
\frac{\opnorm{C}^2\,(\opnorm{D}+E_{\max})}{M\,g_M(E_{\max})}.
\end{equation}
In particular, if $M>\opnorm{D}+E_{\max}$, then
\begin{equation}\label{eq:specA1_crude}
\dist\bigl(E,\spec(A_M^{(1)})\bigr)
\le
\frac{\opnorm{C}^2\,(\opnorm{D}+E_{\max})}{M(M-\opnorm{D}-E_{\max})}.
\end{equation}
\end{theorem}
\proofsuppref{app:proof:thm:first_order_spectral_localization}

\begin{theorem}[Copy-side leakage bound]\label{thm:eigvec_leak}
If $(E,\psi)$ is an eigenpair of $H_{\pm,M}$ with $|E|\le E_{\max}$ and $g_M(E_{\max})>0$,
the copy component satisfies $\|\psi_{\mathrm{copy}}\|\le
\opnorm{C}\,\|\psi_{\mathrm{orig}}\|/g_M(E_{\max})=O(1/M)$.
\end{theorem}
\proofsuppref{app:proof:thm:eigvec_leak}

\begin{theorem}[Inertia reduction for large scalar jamming]\label{thm:inertia_reduction}
Assume $M>\opnorm{D}$, so $D_M=D+M\Id$ is positive, and set $S_M(0):=A-CD_M^{-1}C^\dagger$. Then
$N_-(H_{\pm,M})=N_-(S_M(0))$, $\dim\ker H_{\pm,M}=\dim\ker S_M(0)$, and
$\|S_M(0)-A\|\le\opnorm{C}^2/(M-\opnorm{D})$. Consequently, if $0\notin\spec(A)$ and the
perturbation satisfies $\opnorm{C}^2/(M-\opnorm{D})<\dist(0,\spec(A))$, then $H_{\pm,M}$ is
invertible and $N_-(H_{\pm,M})=N_-(A)$.
\end{theorem}
\proofsuppref{app:proof:thm:inertia_reduction}

\begin{corollary}[Monotonicity in the scalar jam scale]\label{cor:scalar_jam_monotonicity}
Let $M_2>M_1$. The eigenvalues $\lambda_j(H_{\pm,M})$, listed in nondecreasing order with
multiplicity, are nondecreasing in $M$. In particular, $M\mapsto N_-(H_{\pm,M})$ is nonincreasing.
\end{corollary}
\proofsuppref{app:proof:cor:scalar_jam_monotonicity}

\begin{corollary}[No low-energy pure-copy states]\label{cor:no_copy_spectral_subspace}
Let $E_{\max}>0$ and assume $M>\opnorm{D}+E_{\max}$. The compression
$P_{\mathrm{orig}}\,Q_M:\Ran Q_M\to\Hil_{\mathrm{orig}}$ is injective for
$Q_M=\chi_{[-E_{\max},E_{\max}]}(H_{\pm,M})$, so $\rank Q_M\le\dim\Hil_{\mathrm{orig}}$, and
quantitatively $\|P_{\mathrm{copy}}\psi\|\le (E_{\max}+\opnorm C)/(M-\opnorm D)\,\|\psi\|$ for
every $\psi\in\Ran Q_M$.
\end{corollary}
\proofsuppref{app:proof:cor:no_copy_spectral_subspace}

\begin{proposition}[Inherited gap from the original-side block]\label{prop:inherited_gap_A}
Assume $\spec(A)\cap(-g_A,g_A)=\emptyset$ for some $g_A>0$ and
$g_M(g_A):=\dist([-g_A,g_A],\spec(D_M))>0$. With $\Delta_M=\opnorm{C}^2/g_M(g_A)$, if
$\Delta_M<g_A$ then $\spec(H_{\pm,M})\cap(-g_A+\Delta_M,g_A-\Delta_M)=\emptyset$.
\end{proposition}
\proofsuppref{app:proof:prop:inherited_gap_A}

\begin{corollary}[Uniform direct-sum large-jam estimates]\label{cor:uniform_direct_sum_jam}
Under uniform bounds $\bar C=\sup_n\opnorm{C_n}<\infty$ and $\bar D=\sup_n\opnorm{D_n}<\infty$,
the finite-$M$ estimates of
\Cref{thm:resolvent,thm:first_order_eff,thm:spectral_localization,thm:eigvec_leak} hold uniformly
in $n$ for the direct-sum doubled-and-jammed family. The compressed resolvent error, the
first-order compressed approximation, the spectral-localization distance, and the copy leakage
carry the same explicit bounds with $\opnorm{C}$ replaced by $\bar C$ and $\opnorm{D}$ by
$\bar D$.
\end{corollary}
\proofsuppref{app:proof:cor:uniform_direct_sum_jam}

\begin{theorem}[Internal bulk spectral equivalence]\label{thm:internal_bulk_spectral_equivalence}
Let $(\K_n)_n$ be finite triangulated oriented surfaces with full boundary subcomplex, with
$d_n=\dim\Hil(\K_n)$ and uniform bounds on $\|A_{\pm,n}\|,\|H^{\mathrm{PL}}_\pm(\K_n)\|,
\|C_{\pm,n}\|,\|D_{\pm,n}\|$. Let $M_n\to\infty$ and assume $r_n/d_n\to0$ for the
boundary-collar dimension $r_n$ from Theorem~\ref{thm:relative_boundary_comparison}. Then for
every $f\in C_0(\mathbb R)$,
\[
\frac1{d_n}\operatorname{Tr} f(H_{\pm,n,M_n})
-
\frac1{d_n}\operatorname{Tr} f\!\left(H^{\mathrm{PL}}_{\pm}(\K_n)\right)
\longrightarrow0.
\]
The intrinsic Poincar\'e--Lefschetz model, the doubled original-side compression, and the
finite-\(M_n\) jammed model have the same normalized bulk spectral subsequential limits. The
statement uses no closed-bulk Chern theorem, no uniform doubled-bulk gap, and no
interface-pairing hypothesis.
\end{theorem}

\paragraph{Proof outline of \Cref{thm:internal_bulk_spectral_equivalence}.}
The proof combines three building blocks already established in this section; the full version
is in \Cref{S-app:proof:thm:internal_bulk_spectral_equivalence}.
(i) The compressed-resolvent estimate \Cref{thm:resolvent} gives
\(\|P_{\mathrm{orig}}(H_{\pm,n,M_n}-z)^{-1}P_{\mathrm{orig}}-(A_n-z)^{-1}\|\to 0\) as
\(M_n\to\infty\) at the explicit \(O(1/M_n)\) rate; in particular
\(\tfrac{1}{d_n}\bigl|\operatorname{Tr} f(H_{\pm,n,M_n})-\operatorname{Tr} f(A_n)\bigr|\to 0\) for every
\(f\in C_0(\mathbb R)\) by Stone--Weierstrass density of resolvents and the uniform bound
\(d_n^{-1}\operatorname{Tr}|\cdot|\le \|\cdot\|\).
(ii) The collar-rank lemma \Cref{lem:collar_rank_counts}, applied to
\(A_n=H^{\mathrm{PL}}_\pm(\K_n)\) and \(A_n+R_{\partial,\pm,n}=H^{\bd}_{\pm}(\K_n)\), bounds
the normalized trace difference by \(O(r_n/d_n)\to 0\).
(iii) Combining (i) with the identification \(A_n=H^{\bd}_\pm(\K_n)\) and (ii) gives
\(\tfrac{1}{d_n}\bigl|\operatorname{Tr} f(H_{\pm,n,M_n})-\operatorname{Tr} f(H^{\mathrm{PL}}_\pm(\K_n))\bigr|\to 0\) for
every \(f\in C_0(\mathbb R)\). No closed-bulk Chern input or uniform doubled gap is needed; the
hypothesis is only the uniform finite-dimensional operator-norm bound on the blocks
\(A_{\pm,n},C_{\pm,n},D_{\pm,n}\), and the negligible-collar condition \(r_n/d_n\to 0\).

\proofsuppref{app:proof:thm:internal_bulk_spectral_equivalence}

The remaining estimates in this section -- contour eigenvalue-count stability, monotonicity in the
jam scale, no-copy low-energy bounds, an inherited-gap criterion, and the uniform direct-sum
version used by the Roe construction -- are stated together with their explicit constants and
proved in the supplement, Section~\ref{S-sec:schur_proofs}.

%====================================================================
\section{Roe-algebra identification of the doubled-and-jammed interface class}\label{sec:pathb_internalization}
%====================================================================

This section establishes the boundary-map identification used by the headline Roe-winding theorem.
The closed-bulk gap and Chern value come from the closed universal-model theorem of Higson and
Prodan \cite{HigsonProdanUniversalPRL,HigsonProdanUniversalArxiv}. The Ewert--Meyer citation
\cite{EwertMeyer2019} fixes the coarse bulk--interface language and the Mayer--Vietoris pairing
convention; for the doubled-and-jammed original/copy partition used in this paper, the
identification of the Mayer--Vietoris boundary with the Roe connecting map is proved here.

\subsection{Higson--Prodan universal closed-bulk input}

\paragraph{Internal alternative to HP for the flat-torus case.}\label{par:internal_torus_alternative}
The Higson--Prodan closed universal-model theorem \Cref{thm:hp_imported} is invoked below for
\emph{general} bounded-geometry closed triangulations. However, for the special case of a
\emph{flat 2-torus with a uniform triangular triangulation}---a finite quotient
\(T_N=\mathbb R^2/(N\mathbb Z)^2\) equipped with the \(\Lambda\)-periodic triangular
triangulation \(\K_N^\triangle\) of \Cref{sec:periodicbulk}---the closed-bulk Fermi gap of
\Cref{thm:hp_imported}(i) is now proved internally, with explicit constant
\(c=2^{-13}\) independent of the side length \(N\); see
\Cref{S-thm:internal_torus_gap} in the supplement. In particular, the flat-torus instance of the
closed-bulk gap input no longer depends on HP. The Chern integer part of
\Cref{thm:hp_imported}(ii) for the same flat-torus instance is supplied internally at the
\emph{operator level} by \Cref{S-thm:internal_torus_chern} in the supplement
(\(\Ch_N(P_+(\K_N^\triangle))=-1\) and \(\Ch_N(P_-(\K_N^\triangle))=+1\) on the finite torus
\(T_N\) for every \(N\) at least an explicit threshold \(N_0\)), and at the
\emph{symbol level} by \Cref{thm:periodic_triangular_chern} (continuous Chern signs
\(\Ch(E_-(H_+^\triangle))=-1\) and \(\Ch(E_-(H_-^\triangle))=+1\)). The closed-bulk gap input
of \Cref{thm:hp_imported}(i) for the \emph{hexagonal benchmark Bloch model}
\(H_{\hexagon}(k;m)\) of the supplement (\Cref{S-def:hex_bloch_symbol}) on the same flat
2-torus \(T_N\) is also proved internally, at the operator level, by
\Cref{S-thm:internal_hex_gap} (uniform Fermi gap
\(\gamma_{\hexagon}(m)=\sqrt{\min\{1/4,(\sqrt{3}-1)^2m^2\}}\), independent of \(N\)); the
underlying \(2\times 2\) determinant bound is the symbol-level
\Cref{S-lem:hex_bulk_gap}. The companion benchmark---the \emph{square Bloch model}
\(H_{\square}(k;m)\) of the supplement (\Cref{S-def:square_bloch_symbol})---admits the same
internal closed-bulk gap on the flat-torus quotient \(T_N\), proved at the operator level by
\Cref{S-thm:internal_square_gap} (uniform Fermi gap
\(\gamma_{\square}(m)=\sqrt{\min\{1/4,(\sqrt{3}-1)^2m^2\}}\), independent of \(N\)); the
underlying \(2\times 2\) determinant bound is the symbol-level
\Cref{S-lem:square_bulk_gap}. The closed-bulk Chern integer input of
\Cref{thm:hp_imported}(ii) for the \emph{hexagonal benchmark Bloch model} on the same flat
2-torus \(T_N\) is also proved internally at the operator level by
\Cref{S-thm:internal_hex_chern} (\(\Ch_N(P_-(\K_N^{\hexagon};m))=-\sgn(m)\) on the finite torus
\(T_N\) for every \(N\) at least an explicit threshold \(N_0^{\hexagon}(m)\)); the underlying
symbol-level Chern signs come from \Cref{S-lem:hex_chern_positive,S-lem:hex_chern_negative} and
are lifted to the operator level by the same Fukui--Hatsugai--Suzuki convergence mechanism used
for the triangular case in \Cref{S-thm:internal_torus_chern}. The companion benchmark---the
\emph{square Bloch model} \(H_{\square}(k;m)\) of the supplement
(\Cref{S-def:square_bloch_symbol})---admits the same internal closed-bulk Chern integer on the
flat-torus quotient \(T_N\), proved at the operator level by \Cref{S-thm:internal_square_chern}
(\(\Ch_N(P_-(\K_N^{\square};m))=+\sgn(m)\) on the finite torus \(T_N\) for every \(N\) at
least an explicit threshold \(N_0^{\square}(m)\)); the underlying symbol-level Chern sign comes
from \Cref{S-lem:square_chern_positive} (for \(m>0\)) and the antiunitary mass-reversal
argument in the proof of \Cref{S-lem:square_toeplitz_winding} (for \(m<0\)), and is lifted to
the operator level by the same Fukui--Hatsugai--Suzuki convergence mechanism used for the
triangular and hexagonal cases. For flat-torus closed doubles, both the closed-bulk gap
(\Cref{S-thm:internal_torus_gap}) and the closed-bulk Chern integer
(\Cref{S-thm:internal_torus_chern}) are proved internally; both two-band benchmarks of
\Cref{S-app:benchmark_hex_square} (hexagonal and square) on the same flat-torus quotient are
also internal for both inputs (gap: \Cref{S-thm:internal_hex_gap} and
\Cref{S-thm:internal_square_gap}; Chern: \Cref{S-thm:internal_hex_chern} and
\Cref{S-thm:internal_square_chern}); the general bounded-geometry case still cites HP.

\begin{definition}[Universal Hamiltonian on a bounded-geometry triangulation]\label{def:pathb_universal_hamiltonian}
Let \(K\) be a bounded-geometry closed oriented triangulated surface, let
\(\Hil(K)=\ell^2(K_0\sqcup K_1\sqcup K_2)\), let \(B\) be the simplicial boundary operator, and let
\(S_K\) be the Higson--Prodan Hermitian Poincar\'e-duality operator constructed from cap product
with the fundamental cycle. The two universal closed Hamiltonians are
\[
H_\pm(K):=(B+B^\dagger)\pm S_K
\]
on \(\Hil(K)\).
\end{definition}

\begin{theorem}[Cited: Higson--Prodan closed-bulk gap and Chern]\label{thm:hp_imported}
\emph{(Higson--Prodan \cite{HigsonProdanUniversalPRL,HigsonProdanUniversalArxiv}; statement as
imported.)} Let \((K_n)_n\) be a sequence of finite closed oriented triangulated surfaces of
uniformly bounded local geometry: there exist constants \(\Delta\), \(L\), \(R_0\) such that
every vertex of \(K_n\) has simplicial-star valence at most \(\Delta\), every face has at most
\(L\) edges in its star, and the link of every simplex is contained in a combinatorial ball of
radius at most \(R_0\); moreover \((K_n)\) is a refinement sequence in the sense that each
\(K_{n+1}\) admits a simplicial subdivision map to \(K_n\). Let \(H_\pm(K_n)\) be the universal
Chern Hamiltonians of \Cref{def:pathb_universal_hamiltonian}. Then:
\begin{enumerate}[label=(\roman*),leftmargin=2em]
\item (Uniform Fermi gap.) There exists \(g_{\mathrm{HP}}>0\), depending only on
\((\Delta,L,R_0)\) and not on \(n\), such that
\[
\spec\bigl(H_\pm(K_n)\bigr)\cap(-g_{\mathrm{HP}},g_{\mathrm{HP}})=\varnothing
\qquad\text{for all }n.
\]
\item (Closed-bulk Chern integer.) Let
\(P_\pm(K_n)=\chi_{(-\infty,0)}(H_\pm(K_n))\) be the Fermi projection. The closed-bulk
Bellissard--Connes/Higson--Prodan Chern pairing is
\[
\Ch\bigl([P_\pm(K_n)]\bigr)=\pm 1,
\]
with the sign determined by the chosen component \(H_+\) versus \(H_-\) and by the orientation
convention of Section~\ref{sec:setup} (Bloch torus orientation \(dk_1\wedge dk_2\) when the
closed triangulation is periodic, and the canonical Higson--Prodan orientation otherwise).
\end{enumerate}
This is the only external theorem cited anywhere in the proofs of this paper. Bibliographic
data for the published PRL and the arXiv preprint are listed under
\cite{HigsonProdanUniversalPRL} (Phys.\ Rev.\ Lett.\ \textbf{136} (2026), no.~9, 096602,
doi:10.1103/66x5-dvqq) and \cite{HigsonProdanUniversalArxiv} (arXiv:2510.20862,
doi:10.48550/arXiv.2510.20862).
\end{theorem}

\begin{lemma}[Doubled-jammed families satisfy HP hypotheses]\label{lem:hp_hypothesis_check}
Let \((\K_n)_n\) be a uniformly regular boundary refinement family of finite triangulated
oriented surfaces with full boundary subcomplex (in the sense of
Definition~\ref{def:surfwithbdry} together with uniform bounds on local valences and stars),
and let \(D(\K_n)\) be the simplicial doubles of \Cref{prop:double_closed}. Then the closed
sequence \((D(\K_n))_n\) satisfies the hypotheses of \Cref{thm:hp_imported}. Concretely:
\begin{enumerate}[label=(\roman*),leftmargin=2em]
\item (Bounded geometry.) The vertex valences, face-star sizes, and link radii of \(D(\K_n)\)
are bounded uniformly in \(n\) by constants depending only on the corresponding bounds for
\((\K_n)\) and on the doubling operation. Specifically, if every interior vertex of \(\K_n\)
has valence at most \(\Delta_{\mathrm{int}}\) and every boundary vertex has valence at most
\(\Delta_{\bd}\), then every vertex of \(D(\K_n)\) has valence at most
\(\max(\Delta_{\mathrm{int}},2\Delta_{\bd}-1)\); analogous bounds hold for face-star sizes
and link radii by \Cref{prop:double_closed}.
\item (Refinement structure.) If \((\K_n)\) is a refinement sequence (each \(\K_{n+1}\)
subdivides \(\K_n\) and the subdivisions restrict to refinements of \(\bd\K_n\)), then
\((D(\K_n))\) is a refinement sequence in the sense of \Cref{thm:hp_imported}: the doubled
subdivision \(D(\K_{n+1})\) admits a simplicial map to \(D(\K_n)\) extending the subdivision
on each side, gluing continuously across the shared boundary by the full-subcomplex condition
of Definition~\ref{def:surfwithbdry}.
\item (Closed oriented surface.) Each \(D(\K_n)\) is a finite closed oriented triangulated
surface (\Cref{prop:double_closed}); the orientation on the original side is the orientation
of \(\K_n\), and on the copy side is the reversed orientation, consistent across \(\bd\K_n\).
\end{enumerate}
Consequently \Cref{thm:hp_imported} applies to \((D(\K_n))_n\): there is a uniform Fermi gap
\(g_{\mathrm{HP}}>0\) for the closed doubled universal Hamiltonians \(H_\pm(D(\K_n))\), and
their closed-bulk Chern integers are \(\Ch([P_\pm(D(\K_n))])=\pm 1\).
\end{lemma}

\begin{proof}[Proof of \Cref{lem:hp_hypothesis_check}]
The valence, face-star, and link bounds are local combinatorial counts in the doubled complex
\(D(\K_n)=\K_n\cup_{\bd\K_n}\overline{\K_n}\); each interior simplex of \(D(\K_n)\) is also an
interior simplex of either \(\K_n\) or \(\overline{\K_n}\), inheriting its valence/star
bounds, while each boundary simplex of \(\K_n\) is shared with \(\overline{\K_n}\) and has
valence equal to the sum of valences from the two sides. The full-subcomplex condition
\Cref{def:surfwithbdry} ensures that the gluing across \(\bd\K_n\) does not create duplicate
simplices, so \(D(\K_n)\) is an abstract simplicial complex with these uniform bounds. The
refinement property follows from the gluing-compatible subdivision of \(\bd\K_n\) and the
two-sided subdivision of \(\K_n\) and \(\overline{\K_n}\). That \(D(\K_n)\) is a finite
closed oriented triangulated surface is \Cref{prop:double_closed}.
\end{proof}

\begin{lemma}[Higson--Prodan closed-bulk gap, cited]\label{lem:pathb_hp_gap}
For the universal Hamiltonians \(H_\pm(K)\) on closed bounded-geometry triangulations covered
by \Cref{thm:hp_imported}, there is a uniform Fermi gap
\(\spec(H_\pm(K))\cap(-g_{\mathrm{HP}},g_{\mathrm{HP}})=\varnothing\) for some
\(g_{\mathrm{HP}}>0\) independent of the refinement level. (Restatement of
\Cref{thm:hp_imported}(i) for the closed-bulk gap input used throughout
Section~\ref{sec:pathb_internalization}.)
\end{lemma}

\begin{lemma}[Higson--Prodan closed-bulk Chern number, cited]\label{lem:pathb_hp_chern}
With \(P_\pm(K)=\chi_{(-\infty,0)}(H_\pm(K))\), the Bellissard--Connes/Higson--Prodan closed-bulk
Chern pairing satisfies \(\Ch([P_\pm])=\pm1\), with the sign determined by the chosen component
and by the orientation convention of Section~\ref{sec:setup}. (Restatement of
\Cref{thm:hp_imported}(ii) for the closed-bulk Chern input.)
\end{lemma}

\subsection{Roe extension and the interface boundary class}

Let \(X=\bigsqcup_n X(D(\K_n))\) be the coarse disjoint union of the doubled barycentric vertex
sets, let \(X^+\) be the original side, let \(X^-\) be the reflected-copy side, and let \(Y\) be a
fixed coarse collar of the glued boundary. The choice of coarse collar radius does not affect the
ideal $C^*(Y\subset X)$ (Proposition~\ref{S-prop:collar_radius_independence}).

\begin{lemma}[Roe algebra and closed-bulk Fermi projection]\label{lem:pathb_roe_setup}
For \(X\) as above, the scalar Roe algebra \(C^*(X)\) is unital. Under the Higson--Prodan gap of
Lemma~\ref{lem:pathb_hp_gap}, the direct-sum Fermi projection
\(P_\pm=\chi_{(-\infty,0)}\!\left(\bigoplus_n H_\pm(D(\K_n))\right)\) belongs to \(C^*(X)\).
\end{lemma}

\begin{lemma}[Coarse partition exact sequence]\label{lem:pathb_exact_sequence}
With \(\mathcal A=C^*(X)\) and \(\mathcal I=C^*(Y\subset X)\), the original/copy decomposition gives
the quotient extension
\(0\to\mathcal I\to\mathcal A\xrightarrow{\pi}\mathcal A/\mathcal I\to0\), with boundary map
\(\partial_Y:K_0(\mathcal A/\mathcal I)\to K_1(\mathcal I)\).
\end{lemma}

\begin{definition}[Interface projection and unitary]\label{def:pathb_interface_projection}
Let \(J_{\pm,M}^\infty=\bigoplus_n H_{\pm,M,n}\) be the doubled-and-jammed family and let
\(p_J:=E_+\bigl(\pi(P_\pm)\bigr)\in\mathcal A/\mathcal I\). For any admissible cutoff \(h\),
\(p_J=\pi(h(J_{\pm,M}^\infty))\), and the associated interface \(K_1\)-class is represented by
\(u_{J,h}:=\exp(2\pi i\,h(J_{\pm,M}^\infty))\in\Id+\mathcal I\),
\([u_{J,h}]=\partial_Y([p_J])\).
\end{definition}

\subsection{Conventions for Theorem~\ref{thm:pathb_roe_winding}}\label{sec:pathb_conventions}
The headline theorem fixes one set of conventions, used identically in the body proof and the
appendix-level verifications:
\begin{enumerate}[label=(C\arabic*),leftmargin=2.4em]
\item\label{conv:pathb_rep} \emph{Representation.} The standard left-regular
\(\ell^2\)-representation of the scalar Roe algebra: \(C^*(X)\subset\mathcal B(\ell^2(X))\) is the
norm closure of locally compact finite-propagation operators acting on the canonical basis of unit
point masses at each vertex of $X=\bigsqcup_n X(D(\K_n))$.
\item\label{conv:pathb_stab} \emph{Stabilization.} The scalar (not matrix-stabilized) Roe algebra.
Standard Murray--von Neumann \(K\)-theoretic stabilization is invoked only where strictly
necessary, namely in step (3) of the proof below.
\item\label{conv:pathb_sign} \emph{Sign convention for \(\partial\).} The connecting map
\(\partial_Y\) is the standard Higson--Roe boundary map in the six-term exact sequence of
\(0\to\mathcal I\to\mathcal A\to\mathcal A/\mathcal I\to0\), with the orientation choice that
\(\partial[\pi(P)]=[\exp(2\pi i\,\tilde P)]\) for any self-adjoint lift \(\tilde P\) of the
quotient projection \cite{HigsonRoeI,RoeLectures}.
\item\label{conv:pathb_side} \emph{Side convention.} \(X=X^+\sqcup X^-\) is the original/copy
decomposition; spectral cutoffs are taken on the negative real axis,
\(P_\pm=\chi_{(-\infty,0)}(H_\pm^\infty)\).
\item\label{conv:pathb_quotient} \emph{Quotient map.} \(\pi:\mathcal A\to\mathcal A/\mathcal I\) is
the canonical \(C^*\)-quotient by \(\mathcal I=C^*(Y\subset X)\), with collar-radius independence.
\end{enumerate}
The Ewert--Meyer Mayer--Vietoris boundary map for the same original/copy partition is denoted
\(\partial^{\mathrm{EM}}_Y\); the internal compatibility theorem below identifies it with
\(\partial_Y\) under \ref{conv:pathb_rep}--\ref{conv:pathb_quotient}.

\begin{figure}[H]
\centering
\[
\begin{tikzcd}[row sep=2.6em, column sep=2.6em]
0 \arrow[r] & \mathcal I \arrow[r, hook] \arrow[d, equal] & \mathcal A
  \arrow[r, two heads, "\pi"] \arrow[d, equal] & \mathcal A/\mathcal I \arrow[r]
  \arrow[d, equal] & 0 \\
0 \arrow[r] & C^*(Y\subset X) \arrow[r, hook] & C^*(X)
  \arrow[r, two heads, "\pi"] & C^*(X)/C^*(Y\subset X) \arrow[r] & 0 \\
K_1(\mathcal I) \arrow[r, "\cong"] & {[u_{J,h}]}
  & {[P_\pm]\in K_0(\mathcal A)} \arrow[r, "\pi_*"] \arrow[d, "\Ch"', "\mathrm{HP}"]
  & {[p_J]=[E_+(\pi(P_\pm))]\in K_0(\mathcal A/\mathcal I)}
    \arrow[ll, bend right=18, dashed, "\partial_Y\,=\,\partial^{\mathrm{EM}}_Y"'] \\
  & & \mathbb Z\ni\pm1 &
\end{tikzcd}
\]
\caption{Commutative diagram for Theorem~\ref{thm:pathb_roe_winding}. The dashed arrow labelled
\(\partial_Y=\partial^{\mathrm{EM}}_Y\) records the internal compatibility theorem: the Roe
connecting map and the Ewert--Meyer Mayer--Vietoris boundary coincide for the original/copy
partition, sending \([p_J]\) to the interface \(K_1\)-class
\([u_{J,h}]=[\exp(2\pi i\,h(J_{\pm,M}^\infty))]\). The downward arrow \(\mathrm{HP}\) is
Higson--Prodan's closed-bulk Chern character, which evaluates the bulk class to \(\pm1\).}
\label{fig:pathb_diagram}
\end{figure}

\subsection{Internal compatibility theorem}

\begin{theorem}[Roe interface class from the Higson--Prodan closed bulk]\label{thm:pathb_roe_winding}
Let \((\K_n)_n\) be a uniformly regular boundary refinement family whose closed doubles
\(D(\K_n)\) satisfy the hypotheses of \Cref{thm:hp_imported} (verified for the doubled-jammed
families of this paper in \Cref{lem:hp_hypothesis_check}). Use the conventions
\ref{conv:pathb_rep}--\ref{conv:pathb_quotient}. Let \(\mathcal A=C^*(X)\),
\(\mathcal I=C^*(Y\subset X)\), and
\[
P_\pm=\chi_{(-\infty,0)}\!\left(\bigoplus_n H_\pm(D(\K_n))\right).
\]
For every \(M>B_{\mathrm{copy},-}\), the doubled-and-jammed family \(J_{\pm,M}^\infty\)
determines the quotient projection \(p_J=E_+(\pi(P_\pm))\in\mathcal A/\mathcal I\) and the
\(K_1\)-class identity
\[
\Indint(J_{\pm,M}^\infty;T^+,T^-)
=
\partial_Y(E_+(\pi(P_\pm)))
=
[\,\exp(2\pi i\,h(J_{\pm,M}^\infty))\,]\in K_1(\mathcal I)
\]
for any admissible cutoff \(h\). The Roe connecting map \(\partial_Y\) coincides on
\([p_J]\) with the Ewert--Meyer Mayer--Vietoris boundary map \(\partial^{\mathrm{EM}}_Y\)
attached to the original/copy partition under
\ref{conv:pathb_rep}--\ref{conv:pathb_quotient}; this identification is proved internally in
Step (2) of the body proof and is the dashed arrow of Figure~\ref{fig:pathb_diagram}. The
class \(\partial_Y(E_+(\pi(P_\pm)))\) is independent of the admissible cutoff \(h\) and of
the jam scale \(M>B_{\mathrm{copy},-}\), and is preserved under refinement subsequences for
which the Higson--Prodan Fermi gap is uniform.
\end{theorem}

\begin{figure}[H]
\centering
\[
\begin{tikzcd}[row sep=2.4em, column sep=1.6em]
0 \arrow[r] & C^*(Y\subset X) \arrow[r, hook, "\iota"] \arrow[d, equal]
  & C^*(X) \arrow[r, two heads, "\pi"] \arrow[d, equal]
  & C^*(X)/C^*(Y\subset X) \arrow[r] \arrow[d, equal]
  & 0 \\
0 \arrow[r] & \mathcal I \arrow[r, hook, "\iota"]
  & \mathcal A \arrow[r, two heads, "\pi"]
  & \mathcal A/\mathcal I \arrow[r] & 0 \\
& K_1(\mathcal I)
  & K_0(\mathcal A/\mathcal I)
       \arrow[l, dashed, bend left=18, "\partial_Y\,\text{(Roe)}"']
       \arrow[l, dashed, bend right=18, "\partial^{\mathrm{EM}}_Y\,\text{(Ewert--Meyer)}"]
  & K_0(\mathcal A) \arrow[l, "\pi_*"'] \\
& & {[p_J]\,\in\,K_0(\mathcal A/\mathcal I)}
       \arrow[ul, dashed, "\partial_Y\,=\,\partial^{\mathrm{EM}}_Y"', bend left=10]
  &
\end{tikzcd}
\]
\caption{SES-comparison diagram for the body proof of
\Cref{thm:pathb_roe_winding}. \textbf{Top row:} the short exact sequence of \(C^*\)-algebras
\(0\to C^*(Y\subset X)\to C^*(X)\to C^*(X)/C^*(Y\subset X)\to 0\) on which the Higson--Roe
connecting morphism \(\partial_Y\) is built. \textbf{Middle row:} the same SES with the
algebra-symbol abbreviations \(\mathcal I=C^*(Y\subset X)\) and \(\mathcal A=C^*(X)\) of
Convention~\ref{conv:pathb_quotient}; the two rows are identical, displayed as parallel
rows so the connecting-morphism comparison can be read off without renaming. \textbf{Third
row:} the two a priori distinct connecting morphisms attached to this single SES are the
Roe connecting map \(\partial_Y\colon K_0(\mathcal A/\mathcal I)\to K_1(\mathcal I)\) of
Convention~\ref{conv:pathb_sign} (upper bent dashed arrow) and the Ewert--Meyer
Mayer--Vietoris boundary map \(\partial^{\mathrm{EM}}_Y\colon K_0(\mathcal A/\mathcal I)\to
K_1(\mathcal I)\) of \cite{EwertMeyer2019} (lower bent dashed arrow); both have common
source \(K_0(\mathcal A/\mathcal I)\) and common target \(K_1(\mathcal I)\), drawn as two
parallel arrows so the equality is visually compared. The unbent \(\pi_*\) arrow on the
right is the \(K_0\)-image of the quotient map \(\pi\colon\mathcal A\twoheadrightarrow
\mathcal A/\mathcal I\) of the middle row, retained for orientation only.
\textbf{Bottom row:} the cycle representative \([p_J]\in K_0(\mathcal A/\mathcal I)\) on
which the equality is verified; Step (2) of the body proof of
\Cref{thm:pathb_roe_winding} establishes the identity
\(\partial_Y=\partial^{\mathrm{EM}}_Y\) on this class for the original/copy partition under
\ref{conv:pathb_rep}--\ref{conv:pathb_quotient}, and the diagonal dashed arrow displays
this coincidence as a single arrow into \(K_1(\mathcal I)\). The cited literature for the
Mayer--Vietoris boundary in the sense used here is \cite{EwertMeyer2019}; for the Roe
connecting morphism the standard reference is \cite{HigsonRoeI,HigsonRoeII,HigsonRoeIII}. The comparison is the
dashed arrow of \Cref{fig:pathb_diagram} in commutative-diagram form; the present figure
separates the SES structure from the cycle representatives so the parallel-arrow
comparison of the two connecting morphisms can be read in one panel.}
\label{fig:roe_em_ses_comparison}
\end{figure}

\begin{proof}[Body proof of Theorem~\ref{thm:pathb_roe_winding}]
\textbf{Roadmap.} The proof has four steps, in the order displayed by
Figure~\ref{fig:pathb_diagram}. Steps (1), (2), and (4) are internal to this paper. Step (3) is
the only step that uses an external theorem, namely Higson--Prodan
\cite{HigsonProdanUniversalPRL,HigsonProdanUniversalArxiv}, packaged as \Cref{thm:hp_imported}.
The convention block \ref{conv:pathb_rep}--\ref{conv:pathb_quotient} fixes the representation,
stabilization, sign, side, and quotient choices used uniformly across all four steps. Throughout
the proof we work with one fixed short exact sequence of \(C^*\)-algebras
\begin{equation}\label{eq:body_short_exact}
0\longrightarrow \mathcal I\;\overset{\iota}{\hookrightarrow}\;\mathcal A
\;\overset{\pi}{\twoheadrightarrow}\;\mathcal A/\mathcal I\longrightarrow 0,
\qquad
\mathcal A=C^*(X),\quad \mathcal I=C^*(Y\subset X),
\end{equation}
to which we attach \emph{two a priori distinct} six-term boundary maps: the Higson--Roe
connecting map \(\partial_Y\) of Convention~\ref{conv:pathb_sign}, and the Ewert--Meyer
Mayer--Vietoris boundary map \(\partial^{\mathrm{EM}}_Y\) of \cite{EwertMeyer2019}. The
non-trivial content of Step (2) is the identity
\(\partial_Y=\partial^{\mathrm{EM}}_Y\) on \(K_0(\mathcal A/\mathcal I)\) for the original/copy
partition.

\medskip\noindent\textbf{Step (1): The Roe boundary map sends \([p_J]\) to the cutoff unitary
class.}
By Lemma~\ref{lem:pathb_roe_setup}, \(P_\pm\in\mathcal A=C^*(X)\) under the Higson--Prodan gap of
Lemma~\ref{lem:pathb_hp_gap} (which is \Cref{thm:hp_imported}(i)). So the quotient projection
\(p_J:=E_+(\pi(P_\pm))\in\mathcal A/\mathcal I\) is well defined as a self-adjoint idempotent in
the quotient \(C^*\)-algebra \(\mathcal A/\mathcal I\) of \eqref{eq:body_short_exact}. The
functional calculus of the quotient-sidewise statement
(Proposition~\ref{S-prop:quotient_sidewise} in the supplement) provides, for any admissible
cutoff \(h\) (continuous, equal to \(1\) on \([B_{\mathrm{copy},+},\infty)\), to \(0\) on
\((-\infty,-B_{\mathrm{copy},-}]\), and monotone in between), a self-adjoint element
\(h(J_{\pm,M}^\infty)\in\mathcal A\) with \(\pi(h(J_{\pm,M}^\infty))=p_J\); see
Definition~\ref{def:pathb_interface_projection} for the explicit lift. Since
\(h(J_{\pm,M}^\infty)\) is a self-adjoint lift of an idempotent in the quotient,
\((h(J_{\pm,M}^\infty))^2-h(J_{\pm,M}^\infty)\in\mathcal I\), hence
\(\exp(2\pi i\,h(J_{\pm,M}^\infty))\in\Id+\mathcal I\) and therefore defines a class in
\(K_1(\mathcal I)\). Convention~\ref{conv:pathb_sign} fixes the sign of the standard
Higson--Roe \(K\)-theoretic connecting map for the short exact sequence
\eqref{eq:body_short_exact}, and the boundary-map lemma
(Lemma~\ref{S-lem:pathb_boundary_map}) gives the labelled identity
\begin{equation}\label{eq:body_step1}
\partial_Y([p_J])=[\exp(2\pi i\,h(J_{\pm,M}^\infty))]\in K_1(\mathcal I),
\end{equation}
independent of the admissible cutoff \(h\). The class \([\exp(2\pi i\,h(J_{\pm,M}^\infty))]\)
is the leftmost downstream class along the top row of Figure~\ref{fig:pathb_diagram}, and
\eqref{eq:body_step1} is precisely the left dashed arrow.
If the sign convention~\ref{conv:pathb_sign} were replaced by its opposite (the alternate
\(\partial[\pi(P)]=-[\exp(2\pi i\,\tilde P)]\) sometimes seen in the literature), the
right-hand side of \eqref{eq:body_step1} would acquire an overall minus sign, producing the
inverse class in \(K_1(\mathcal I)\) and inverting all integer winding evaluations downstream;
this is the auditable consequence of fixing the sign once and for all.

\medskip\noindent\textbf{Step (2): The Roe and Ewert--Meyer boundary maps coincide for the
original/copy partition.}
This is the key compatibility step. We claim that the Higson--Roe connecting map \(\partial_Y\)
of Convention~\ref{conv:pathb_sign} and the Ewert--Meyer Mayer--Vietoris boundary map
\(\partial^{\mathrm{EM}}_Y\) of \cite{EwertMeyer2019} agree on
\(K_0(\mathcal A/\mathcal I)\); concretely, viewed as homomorphisms
\(K_0(\mathcal A/\mathcal I)\to K_1(\mathcal I)\), they are the \emph{same} homomorphism.

\smallskip
\emph{(2a) Both maps come from one and the same short exact sequence.} The Roe-extension map
\(\partial_Y\) is the boundary map of the six-term exact sequence attached to
\eqref{eq:body_short_exact} with the three named ingredients
\((\mathcal A,\mathcal I,\pi)=(C^*(X),C^*(Y\subset X),\pi_{\mathrm{Roe}})\). The Ewert--Meyer
Mayer--Vietoris boundary \(\partial^{\mathrm{EM}}_Y\) is the boundary map of the six-term exact
sequence attached to the bulk--interface pair
\((X^+,X^-;Y)\) used in \cite{EwertMeyer2019}; for the doubled-and-jammed setting at hand, the
constituent algebras are: the bulk algebra is \(C^*(X)\) for the same \(X=X^+\sqcup X^-\); the
interface ideal is \(C^*(Y\subset X)\) for the same \(Y\); and the quotient is the canonical
\(C^*\)-quotient \(C^*(X)/C^*(Y\subset X)\), with the same \(C^*\)-quotient map. The
coarse-excision statement (Proposition~\ref{S-prop:coarse_excision}) is exactly the assertion
that these triples match: the bulk algebra \(\mathcal A_{\mathrm{EM}}\) of \cite{EwertMeyer2019}
is \(C^*(X)\), the interface ideal \(\mathcal I_{\mathrm{EM}}\) is \(C^*(Y\subset X)\), and
the EM quotient \(\mathcal A_{\mathrm{EM}}/\mathcal I_{\mathrm{EM}}\) is the canonical
\(C^*\)-quotient. We record this comparison as the labelled identity
\begin{equation}\label{eq:body_step2_algebras}
(\mathcal A,\mathcal I,\mathcal A/\mathcal I,\pi)_{\mathrm{Roe}}
\;=\;
(\mathcal A,\mathcal I,\mathcal A/\mathcal I,\pi)_{\mathrm{EM}}.
\end{equation}

\smallskip
\emph{(2b) Both maps are induced by the same lift.} For the Roe extension, the boundary class
\(\partial_Y([\pi(P)])\) is by definition (Convention~\ref{conv:pathb_sign}) the class
\([\exp(2\pi i\,\tilde P)]\) for any self-adjoint lift \(\tilde P\) of \(\pi(P)\); see
\cite[\S6]{HigsonRoeI} and \cite{RoeLectures}. For the Ewert--Meyer extension, the
Mayer--Vietoris boundary on a coarsely excisive partition is the connecting map of the same
six-term exact sequence, with the same lift recipe: a self-adjoint operator on \(\ell^2(X)\) of
finite propagation that maps to \(\pi(P)\) under the quotient is the EM lift; see
\cite[\S3]{EwertMeyer2019} where the Mayer--Vietoris boundary is constructed as the boundary
map of the underlying \(C^*\)-extension. Since the only datum required to define either
boundary map on a single element of \(K_0(\mathcal A/\mathcal I)\) is the lift and the sign
convention, and both are fixed identically by \eqref{eq:body_short_exact} and
Convention~\ref{conv:pathb_sign}, we record the labelled identity
\begin{equation}\label{eq:body_step2_lift}
\partial_Y([\pi(P)])\;=\;[\exp(2\pi i\,\tilde P)]\;=\;\partial^{\mathrm{EM}}_Y([\pi(P)])
\qquad\text{for every self-adjoint lift }\tilde P\text{ of }\pi(P).
\end{equation}
By naturality of the six-term exact sequence (a standard property of any short exact sequence
of \(C^*\)-algebras; \cite[\S 8]{BlackadarKTheory}), two boundary maps for the same short exact
sequence with the same lift and the same sign convention coincide as homomorphisms on
\(K_0(\mathcal A/\mathcal I)\): there is no degree of freedom remaining.

\smallskip
\emph{(2c) Side and quotient consistency.} Convention~\ref{conv:pathb_side} fixes
\(X=X^+\sqcup X^-\) as the original/copy partition; the spectral cutoff is on
\((-\infty,0)\), so \(P_\pm=\chi_{(-\infty,0)}(H_\pm^\infty)\) is the negative Fermi projection.
This is the same side and spectral choice used by \cite{EwertMeyer2019} for the negative-energy
Fermi-projection class; the alternative \(\chi_{(0,\infty)}\) cutoff would produce \(\Id-P_\pm\)
instead, and would flip the sign of every quotient class downstream. The quotient
projection \(p_J=E_+(\pi(P_\pm))\) is computed by the quotient-sidewise identity
(Proposition~\ref{S-prop:quotient_sidewise}); the splitting identity
\(\mathcal A/\mathcal I=E_+(\mathcal A/\mathcal I)\oplus E_-(\mathcal A/\mathcal I)\)
(Proposition~\ref{S-prop:quotient_split}) reduces the comparison to the positive side because
the negative-side comparison operator has trivial Fermi projection by the trivial-jam
statement (Proposition~\ref{S-prop:trivial_jam}). If the side convention~\ref{conv:pathb_side}
were replaced by the opposite partition (copy on the original side), the labelled
identity~\eqref{eq:body_step2_lift} would still hold but with the roles of
\((X^+,X^-)\) interchanged; the trivial-jam statement and the splitting identity guarantee that
the answer does not depend on which side is chosen as \(X^+\), only that the splitting itself
is fixed.

\smallskip
\emph{(2d) Quotient convention.} Both boundary maps use the same \(C^*\)-quotient
\(\pi:\mathcal A\to\mathcal A/\mathcal I\) (Convention~\ref{conv:pathb_quotient}), with
collar-radius independence (Proposition~\ref{S-prop:collar_radius_independence}). An
alternative quotient by a smaller or larger ideal would change \(\mathcal A/\mathcal I\) and
hence both boundary maps; the equality
\(\partial_Y=\partial^{\mathrm{EM}}_Y\) is a statement about the \emph{specific} quotient by
\(C^*(Y\subset X)\), as fixed in \eqref{eq:body_short_exact}, and not about boundary maps for
unrelated ideals.

\smallskip
Combining \eqref{eq:body_step2_algebras}, \eqref{eq:body_step2_lift}, and the side/quotient
checks (2c)--(2d), we conclude the labelled compatibility identity
\begin{equation}\label{eq:body_step2_main}
\partial_Y\;=\;\partial^{\mathrm{EM}}_Y
\qquad\text{on}\qquad K_0(\mathcal A/\mathcal I),
\end{equation}
which is the dashed arrow of Figure~\ref{fig:pathb_diagram} labelled
\(\partial_Y=\partial^{\mathrm{EM}}_Y\). In particular,
\(\partial^{\mathrm{EM}}_Y([p_J])=\partial_Y([p_J])
=[\exp(2\pi i\,h(J_{\pm,M}^\infty))]\) by composing \eqref{eq:body_step1} with
\eqref{eq:body_step2_main}. We emphasise that this identity uses no part of the closed-bulk
Chern statement; only the gap of \Cref{thm:hp_imported}(i) is invoked, via
Lemma~\ref{lem:pathb_roe_setup}, to ensure \(P_\pm\in\mathcal A\). If either the EM lift recipe
or the EM sign convention differed from~\ref{conv:pathb_sign}, the right-hand side of
\eqref{eq:body_step2_lift} could pick up a sign, producing the inverse class in
\(K_1(\mathcal I)\) and breaking the labelled equality \eqref{eq:body_step2_main}; both
references \cite{EwertMeyer2019,HigsonRoeI} use the same six-term-exact-sequence sign and lift
recipe (the difference between the two reduces to a choice of book-keeping for the
Mayer--Vietoris pairing on coarsely excisive partitions), so the convention agreement is
verified.

\medskip\noindent\textbf{Step (3): Closed-bulk Chern from Higson--Prodan.}
By Lemma~\ref{lem:pathb_hp_chern}, which restates \Cref{thm:hp_imported}(ii) for closed
bounded-geometry triangulated surfaces, the closed-bulk Fermi projection has Chern number
\(\Ch([P_\pm])=\pm1\). This is the only step that uses an external theorem, namely
\Cref{thm:hp_imported}; the doubled families satisfy the HP hypotheses by
\Cref{lem:hp_hypothesis_check}. The integer evaluation of the \(K_1\)-class identity from
Step (2) is taken up by the two corollaries (\Cref{cor:toeplitz_winding} for the internal
Toeplitz periodic case; \Cref{cor:conditional_winding} for the general conditional case).
This is the downward Chern-character arrow of Figure~\ref{fig:pathb_diagram} labelled
\(\mathrm{HP}\). If the orientation convention of Section~\ref{sec:setup} were flipped (e.g.\
\(dk_2\wedge dk_1\) instead of \(dk_1\wedge dk_2\) on the Bloch torus), the closed-bulk Chern
sign would flip and the integer winding downstream would change sign correspondingly; the
sign-ledger appendix (Section~\ref{S-app:sign_ledger}) records the explicit chain.

\medskip\noindent\textbf{Step (4): Finite-\(M\) quotient identification and \(K_1\)-limit.}
For every \(M>B_{\mathrm{copy},-}\), the trivial-jam statement
(Proposition~\ref{S-prop:trivial_jam}) gives strict positivity of the negative comparison operator,
so its negative Fermi projection vanishes. The positive comparison operator is
\(T^+=H_\pm^\infty\) with negative Fermi projection \(P_\pm\), and the quotient-sidewise
identification (Proposition~\ref{S-prop:quotient_sidewise}) gives
\(p_J=E_+(\pi(P_\pm))+E_-(\pi(0))=E_+(\pi(P_\pm))\) at every finite jam scale
\(M>B_{\mathrm{copy},-}\). The corresponding \(K_1\)-class is independent of \(M\)
(Corollary~\ref{S-cor:jam_independence}), so the limit \(M\to\infty\) is well defined as a
single \(K_1(\mathcal I)\)-class. Composing the four steps---namely
\eqref{eq:body_step1} from Step (1), \eqref{eq:body_step2_main} from Step (2),
\Cref{thm:hp_imported}(ii) from Step (3), and the quotient-sidewise identification of this
Step (4)---gives the claimed \(K_1\)-class identity
\(\Indint(J_{\pm,M}^\infty;T^+,T^-)=\partial_Y(E_+(\pi(P_\pm)))
=[\exp(2\pi i\,h(J_{\pm,M}^\infty))]\). This is the full top row of
Figure~\ref{fig:pathb_diagram} traversed from
\([P_\pm]\) on the right, through \(\pi_*\) to \([p_J]\), and through the dashed
\(\partial_Y=\partial^{\mathrm{EM}}_Y\) arrow back to \([u_{J,h}]\) on the left. Detailed
computational verifications of the cutoff-functional-calculus identity, the norm-continuous
homotopy along admissible cutoffs, and the explicit lift \(h(J_{\pm,M}^\infty)\) are
collected in the supplement, Section~\ref{S-app:proof:thm:pathb_roe_winding}. If the
quotient-sidewise identification of (4) used a different (e.g.\ Schur-complement) lift, the
\(K_1\)-class \([\exp(2\pi i\,h(J_{\pm,M}^\infty))]\) would still equal
\(\partial_Y([p_J])\) by \eqref{eq:body_step1}, but the explicit lift formula on the
right-hand side would differ; the cutoff-independence of the class
(Corollary~\ref{S-cor:jam_independence}) is what makes the identity robust under the choice
of admissible \(h\).
\end{proof}

\begin{corollary}[Integer winding in the Toeplitz periodic case]\label{cor:toeplitz_winding}
Suppose \(Y\subset X\) is the straight periodic interface of the periodic triangular class of
Section~\ref{sec:periodicbulk}, with primitive tangential/positive-normal basis \((t,n)\) and
\(\varepsilon(t,n)=\det[t\ n]\in\{\pm1\}\). Then the \(K_1\)-class
\(\partial_Y(E_+(\pi(P_\pm)))\) of Theorem~\ref{thm:pathb_roe_winding} maps under the
Toeplitz extension to the integer
\[
\operatorname{wind}\bigl(\partial_Y(E_+(\pi(P_\pm)))\bigr)
=\varepsilon(t,n)\,\Ch\bigl(E_-(H_\pm^\triangle)\bigr)
=\mp\,\varepsilon(t,n)
=\mp\det[t\ n]\in\{\pm 1\},
\]
identifying the \(K_1\)-class with \(\mathbb Z\) via the Toeplitz quotient in
Theorem~\ref{thm:periodic_toeplitz_pairing}. The two displayed equalities use, in order, the
Toeplitz pairing formula \(\operatorname{wind}(\partial_{\mathrm T,t,n}[P])
=\varepsilon(t,n)\,\Ch(P)\) of Corollary~\ref{cor:primitive_periodic_toeplitz} and the Chern
signs \(\Ch(E_-(H_+^\triangle))=-1\), \(\Ch(E_-(H_-^\triangle))=+1\) of
Theorem~\ref{thm:periodic_triangular_chern}; the resulting upper/lower sign reading agrees with
the straight-edge formula \(\operatorname{wind}=-\det[t\ n]\) for \(H_+^\triangle\) and
\(\operatorname{wind}=+\det[t\ n]\) for \(H_-^\triangle\) of
Corollary~\ref{cor:periodic_straight_edge_winding}. The integer-pairing statement in this case
is fully internal: the closed-bulk gap and Chern signs are proved in
Theorem~\ref{thm:periodic_triangular_bulk} and Theorem~\ref{thm:periodic_triangular_chern},
and the Toeplitz identification of \(K_1\) with \(\mathbb Z\) is the standard BDF--Toeplitz
extension recalled in Theorem~\ref{thm:periodic_toeplitz_pairing} and
Corollary~\ref{cor:primitive_periodic_toeplitz}.
\end{corollary}

\paragraph{Proof outline of \Cref{cor:toeplitz_winding}.}
The \(K_1\)-class identity \(\partial_Y(E_+(\pi(P_\pm)))=[\exp(2\pi i\,h(J_{\pm,M}^\infty))]\)
of \Cref{thm:pathb_roe_winding} is the input. Compressing the positive-normal direction \(n\)
of the periodic interface to a half-space realizes the interface ideal
\(\mathcal I=C^*(Y\subset X)\) as the kernel of the Toeplitz extension
\(0\to C(S^1,\mathcal K)\to\mathcal T_{\mathrm{per}}\to C(\mathbb T^2)\to 0\) of
\Cref{thm:periodic_toeplitz_pairing}, whose half-space-direction connecting map
\(\partial_{\mathrm T,t,n}\) satisfies
\(\operatorname{wind}(\partial_{\mathrm T,t,n}[P])=\varepsilon(t,n)\,\Ch(P)\) by
\Cref{cor:primitive_periodic_toeplitz}. With \(P=E_-(H_\pm^\triangle)\) and the Chern
signs of \Cref{thm:periodic_triangular_chern}, the displayed equality holds: the trivial-jam
statement gives \([P_-]=[0]\), so only the positive-side contribution remains. Full
verification of the cutoff-functional-calculus identification of the \(K_1\)-class with the
Toeplitz winding integer is in \Cref{S-app:proof:cor:toeplitz_winding}.

\proofsuppref{app:proof:cor:toeplitz_winding}

\begin{corollary}[Conditional integer winding for general partitions]\label{cor:conditional_winding}
Suppose that the \(K_1\)-class of the interface ideal \(\mathcal I=C^*(Y\subset X)\) admits
an integer-valued pairing \(\langle\,\cdot\,,\,\cdot\,\rangle:K_0(\mathcal A/\mathcal I)\times
K_1(\mathcal I)\to\mathbb Z\) of Chern-jump type, in the sense that
\(\langle [p_J],\partial_Y([p_J])\rangle=\Ch([P_+])-\Ch([P_-])\) whenever both sides are
defined and \(\partial_Y\) is the Higson--Roe connecting map of
Convention~\ref{conv:pathb_sign}. Then the pairing of
\(\partial_Y(E_+(\pi(P_\pm)))\) with \([p_J]\) is \(\pm1\). The Toeplitz periodic case
(Corollary~\ref{cor:toeplitz_winding}) is the special case where the integer pairing is
internal; the general case is conditional on the existence of an Ewert--Meyer/Chern-jump
compatible integer pairing whose construction outside the periodic class is not carried out in
this paper.
\end{corollary}
\proofsuppref{app:proof:cor:conditional_winding}

\begin{corollary}[Finite-\(M\) cutoff stability]\label{cor:pathb_finite_m_winding}
For every \(M>B_{\mathrm{copy},-}\), the \(K_1(\mathcal I)\)-class
\(\partial_Y(E_+(\pi(P_\pm)))\) of Theorem~\ref{thm:pathb_roe_winding} is independent of the
jam scale \(M\): equality of classes in \(K_1(\mathcal I)\) follows from the cutoff
homotopy and jam-scale-independence statements
(Corollary~\ref{S-cor:jam_independence} and Proposition~\ref{S-prop:quotient_sidewise}).
\end{corollary}

\subsection{Stability and cross-checks}

\begin{lemma}[Collar-supported stability without norm-smallness]\label{lem:pathb_collar_stability}
Let $V$ be any bounded self-adjoint perturbation in $C^*(Y\subset X)$, with no small-norm
assumption, and suppose the positive and negative comparison gaps used to define the sidewise
interface remain open. Then the quotient projection and \(K_1\)-class
\(\partial_Y(E_+(\pi(P_\pm)))\) of Theorem~\ref{thm:pathb_roe_winding} are unchanged. In the
straight periodic Toeplitz case (Corollary~\ref{cor:toeplitz_winding}), the resulting integer
winding is also unchanged.
\end{lemma}
\proofsuppref{app:proof:lem:pathb_collar_stability}

\begin{lemma}[Bounded-geometry stability]\label{lem:pathb_bg_stability}
Along a norm-continuous bounded-geometry deformation of the closed doubled universal Hamiltonians
that preserves the Higson--Prodan Fermi gap, the interface \(K_1\)-class of
Theorem~\ref{thm:pathb_roe_winding} is constant.
\end{lemma}
\proofsuppref{app:proof:lem:pathb_bg_stability}

\begin{lemma}[Refinement stability]\label{lem:pathb_refinement_stability}
For a refinement family $K_n\to K_\infty$ in bounded geometry, if the Higson--Prodan theorem gives
a uniform gap $g_{\mathrm{HP}}$ and the jam scales satisfy $M_n>B_{\mathrm{copy},-}+1$,
$M_n\to\infty$, then the quotient class and \(K_1\)-class of
Theorem~\ref{thm:pathb_roe_winding} are eventually constant under subsequences or cofinal
refinements.
\end{lemma}
\proofsuppref{app:proof:lem:pathb_refinement_stability}

\begin{remark}[Periodic triangular cross-check]\label{lem:pathb_periodic_triangular}
For periodic triangular doubles, Theorem~\ref{thm:pathb_roe_winding} agrees with the internal
Bloch gap, Chern, and straight Toeplitz winding calculations of
Section~\ref{sec:periodicbulk}; details in the supplement,
Section~\ref{S-app:proof:lem:pathb_periodic_triangular}.
\end{remark}

\begin{remark}[Finite-defect cross-check]\label{lem:pathb_finite_defect}
For doubled finite-defect families whose defects remain in a fixed collar or in finitely many
coarse components, Theorem~\ref{thm:pathb_roe_winding} gives the same interface class
as the finite-defect robustness statements in the supplement,
\Cref{S-app:robustness_finite_defects}; details in
Section~\ref{S-app:proof:lem:pathb_finite_defect}.
\end{remark}

\begin{remark}[Hexagonal and square cross-check]\label{lem:pathb_hex_square}
For the supplement's hexagonal and square symbols (\Cref{S-app:benchmark_hex_square}), whenever
their closed-bulk class is identified with the corresponding Higson--Prodan closed-bulk class or
with the internally computed periodic Chern class through a gapped homotopy,
Theorem~\ref{thm:pathb_roe_winding} returns the same interface \(K_1\)-class, and the
Toeplitz integer winding agrees with Corollary~\ref{cor:toeplitz_winding} in the periodic
case; details in Section~\ref{S-app:proof:lem:pathb_hex_square}.
\end{remark}

\begin{remark}[Tail-split and annular cross-check]\label{lem:pathb_tail_annular}
For the tail-split and annular families of the supplement
(\Cref{S-app:robustness_tail_annular}), the boundary class of
Theorem~\ref{thm:pathb_roe_winding} restricts componentwise to the same Toeplitz quotient
coordinates used there; details in Section~\ref{S-app:proof:lem:pathb_tail_annular}.
\end{remark}

\begin{remark}[Input synthesis]\label{rem:pathb_synthesis}
A naive accounting separates the closed doubled-bulk gap, the closed-bulk Chern pairing, and the
Ewert--Meyer comparison into three distinct inputs. Theorem~\ref{thm:pathb_roe_winding} packages
the first two as the single Higson--Prodan closed theorem and proves the third for the
doubled-jammed original/copy partition. The only cited closed-model theorem needed for the
\(K_1\)-class identity is the Higson--Prodan theorem. The integer winding evaluation is
internal in the periodic triangular Toeplitz case (\Cref{cor:toeplitz_winding}) and is
recorded as a conditional corollary for general partitions
(\Cref{cor:conditional_winding}).
\end{remark}

The companion coarse interface class construction, which places the doubled-and-jammed refinement
family in a Roe-algebra construction producing the same quotient projection and \(K_1\)-class
from uniform locality, boundedness, and an available closed-bulk gap, is recorded in the
supplement, Section~\ref{S-sec:bulkboundary}. Its main statement
(Theorem~\ref{S-thm:bulkboundary_main}) and the Ewert--Meyer-language restatement
(Theorem~\ref{S-thm:full_hp_bec}) consolidate the same conclusion.

%====================================================================
\section{A periodic triangular bulk input}\label{sec:periodicbulk}
%====================================================================

One closed-bulk input is checked directly here, rather than imported from Higson--Prodan. Let
\(z_j=e^{ik_j}\), \(k=(k_1,k_2)\in\mathbb T^2\), and consider the
\(\mathbb Z^2\)-periodic triangular triangulation with one vertex \(v\), three oriented edges
\(a,b,c\), and two oriented faces \(f,g\) per fundamental cell. The Bloch fiber is
\[
\mathcal H(k)
=
\mathbb C\langle v\rangle\oplus
\mathbb C\langle a,b,c\rangle\oplus
\mathbb C\langle f,g\rangle .
\]
With respect to this ordered basis, the boundary blocks are
\[
\partial_1(k)=
\begin{bmatrix}
z_1-1&z_2-1&z_1z_2-1
\end{bmatrix},
\qquad
\partial_2(k)=
\begin{bmatrix}
1&-z_2\\
z_1&-1\\
-1&1
\end{bmatrix}.
\]
A direct multiplication gives \(\partial_1(k)\partial_2(k)=0\). The cap-product blocks are
\[
P_0(k)=
\begin{bmatrix}
z_1^{-1}z_2^{-1}\\
z_2^{-1}
\end{bmatrix},
\qquad
P_1(k)=
\begin{bmatrix}
0&z_1^{-1}&0\\
0&0&0\\
-z_2^{-1}&0&0
\end{bmatrix},
\qquad
P_2(k)=
\begin{bmatrix}
1&1
\end{bmatrix},
\]
and the Hermitian duality block is
\[
S_0(k)=\tfrac{i}{2}(P_2(k)^*+P_0(k)),\quad
S_1(k)=\tfrac{i}{2}(P_1(k)^*-P_1(k)),\quad
S_2(k)=-\tfrac{i}{2}(P_0(k)^*+P_2(k)),
\]
so that \(H_\pm^\triangle(k)=B(k)+B(k)^*\pm S(k)\) is Hermitian.

\begin{proposition}[Periodic triangular characteristic polynomial]\label{prop:periodic_triangular_charpoly}
Set
\[
F(k)=28-8\cos k_1-6\cos k_2-6\cos(k_1+k_2),
\qquad
G(k)=\sin k_2+\sin(k_1+k_2).
\]
Then
\[
\det\!\left(\lambda\Id-H_\pm^\triangle(k)\right)
=
\bigl(\lambda^2-\tfrac{F(k)}4\bigr)
\bigl(\lambda^4-\tfrac{13}{2}\lambda^2+\tfrac{F(k)}4
\pm \lambda G(k)\bigr),
\]
with one exact pair of eigenvalue branches \(\lambda=\pm \tfrac12\sqrt{F(k)}\).
\end{proposition}
\proofsuppref{app:proof:prop:periodic_triangular_charpoly}

\begin{theorem}[Periodic triangular bulk gap]\label{thm:periodic_triangular_bulk}
For every \(k\in\mathbb T^2\), \(H_\pm^\triangle(k)\) is invertible. More precisely,
\[
\det H_\pm^\triangle(k)
=
-\tfrac1{16}
\left(
-28+8\cos k_1+6\cos k_2+6\cos(k_1+k_2)
\right)^2
\le -4.
\]
Moreover \(\|H_\pm^\triangle(k)\|\le8\), hence
\(\dist\bigl(0,\spec H_\pm^\triangle(k)\bigr)\ge 2^{-13}\) for all \(k\in\mathbb T^2\), and the
negative Bloch bundle has rank \(3\).
\end{theorem}
\proofsuppref{app:proof:thm:periodic_triangular_bulk}

\begin{corollary}[Small periodic perturbations preserve the triangular gap]\label{cor:periodic_small_perturbation_gap}
If \(V(k)\) is a continuous self-adjoint \(6\times6\) matrix over \(\mathbb T^2\) with
\(\sup_k\opnorm{V(k)}<2^{-13}\), then \(H_\pm^\triangle(k)+V(k)\) is invertible for every
\(k\in\mathbb T^2\), and its negative Bloch bundle has rank \(3\).
\end{corollary}
\proofsuppref{app:proof:cor:periodic_small_perturbation_gap}

\begin{proposition}[Positive duality strength homotopy]\label{prop:periodic_positive_mass_homotopy}
For \(m>0\), the operator \(H_{\pm,m}^\triangle(k)=B(k)+B(k)^*\pm mS(k)\) is invertible for every
\(k\in\mathbb T^2\); explicitly,
\[
\det H_{\pm,m}^\triangle(k)
=
-\tfrac1{16}
\bigl(8(3-c_1-c_2-c_{12})+2m^2(2+c_2+c_{12})\bigr)^2,
\]
with \(c_j=\cos k_j\), \(c_{12}=\cos(k_1+k_2)\). All positive values of the duality strength are
connected by a zero-gapped norm-continuous homotopy.
\end{proposition}
\proofsuppref{app:proof:prop:periodic_positive_mass_homotopy}

\begin{theorem}[Periodic triangular Chern signs]\label{thm:periodic_triangular_chern}
For the periodic triangular Bloch Hamiltonians above, the rank-three negative Bloch bundles have
\[
\Ch\bigl(E_-(H_+^\triangle)\bigr)=-1,
\qquad
\Ch\bigl(E_-(H_-^\triangle)\bigr)=+1.
\]
\end{theorem}
\proofsuppref{app:proof:thm:periodic_triangular_chern}

\begin{corollary}[Periodic Chern stability]\label{cor:periodic_chern_stability}
With \(g_{\pm,m}:=\min_k\dist(0,\spec H_{\pm,m}^\triangle(k))>0\), every continuous self-adjoint
periodic perturbation \(V(k)\) with \(\sup_k\opnorm{V(k)}<g_{\pm,m}\) keeps
\(H_{\pm,m}^\triangle(k)+V(k)\) gapped at zero and preserves the Chern number. The same holds for
\(H_\pm^\triangle\) when \(\sup_k\opnorm{V(k)}<2^{-13}\). In particular, sufficiently small
finite-range self-adjoint periodic perturbations of the triangular model do not change the Chern
signs.
\end{corollary}
\proofsuppref{app:proof:cor:periodic_chern_stability}

As a numerical check, the companion script applies the gauge-invariant lattice formula of
Fukui--Hatsugai--Suzuki \cite{FukuiHatsugaiSuzuki} to the rank-three negative spectral projection,
checking the returned integer, the link denominators, and the plaquette branch margins; the
computation records the same signs as Theorem~\ref{thm:periodic_triangular_chern}. The
antiunitary equivalence
\(H_-^\triangle(k)=\overline{H_+^\triangle(-k)}\) explains the opposite signs through complex
conjugation of the Bloch bundle, with the base map \(k\mapsto-k\) orientation-preserving in two
dimensions.

\begin{theorem}[Toeplitz evaluation for a straight periodic interface]\label{thm:periodic_toeplitz_pairing}
Let \(H(k_1,k_2)\) be a continuous finite-range self-adjoint \(N\times N\) Bloch Hamiltonian on
\(\mathbb T^2\), gapped at zero, and let
\(P(k_1,k_2)=\chi_{(-\infty,0)}(H(k_1,k_2))\). Write
\(\Ch(P)=\frac{1}{2\pi i}\int_{\mathbb T^2}\tr(P\,dP\wedge dP)\in\mathbb Z\) with the
orientation \(dk_1\wedge dk_2\). Compress the second lattice direction to a half-space. After
Fourier transform in the first direction, the straight-edge extension is the Toeplitz extension
\(0\to C(S^1,\mathcal K)\to\mathcal T_{\mathrm{per}}\xrightarrow{\sigma}C(\mathbb T^2)\to 0\) with
connecting map \(\partial_{\mathrm T}:K_0(C(\mathbb T^2))\to K_1(C(S^1))\) satisfying
\(\operatorname{wind}(\partial_{\mathrm T}([P]))=\Ch(P)\). The sign is fixed by the above
orientation and by taking the compressed coordinate to be the positive normal direction; see, for
example, \cite{BDF1977,BlackadarKTheory,ProdanSchulzBaldesBook,BourneKellendonkRennie}.
\end{theorem}
\proofsuppref{app:proof:thm:periodic_toeplitz_pairing}

\begin{corollary}[Primitive-direction straight periodic interfaces]\label{cor:primitive_periodic_toeplitz}
Let \(t,n\in\mathbb Z^2\) be a lattice basis with \(t\) the tangential direction of a straight
periodic interface and \(n\) the compressed positive normal direction. With
\(\varepsilon(t,n):=\det[t\ n]\in\{\pm1\}\), the half-space Toeplitz boundary map in the
\(\eta\)-direction (\(k=\kappa t+\eta n\)) satisfies
\(\operatorname{wind}\!\left(\partial_{\mathrm T,t,n}([P])\right)=\varepsilon(t,n)\,\Ch(P)\). In
particular, every primitive straight periodic edge direction is covered by a lattice change of
variables.
\end{corollary}
\proofsuppref{app:proof:cor:primitive_periodic_toeplitz}

\begin{corollary}[Periodic straight-edge winding from Toeplitz pairing]\label{cor:periodic_straight_edge_winding}
For the periodic triangular bulk models \(H^\triangle_\pm\), every primitive straight periodic
interface with tangential/positive-normal basis \((t,n)\) has interface winding
\[
\operatorname{wind}=-\det[t\ n]\quad\text{for }H^\triangle_+,
\qquad
\operatorname{wind}=\det[t\ n]\quad\text{for }H^\triangle_-.
\]
The same formula holds for every fixed positive-duality-strength triangular model and for every
small finite-range periodic perturbation covered by
Corollary~\ref{cor:periodic_chern_stability}. Finite-cover, finite-defect, finite-corner, split,
and tail-split variants are the appendix statements collected in the supplement,
Section~\ref{S-app:section5_variants}.
\end{corollary}
\proofsuppref{app:proof:cor:periodic_straight_edge_winding}

Hexagonal and square two-band benchmarks are recorded in the supplement,
Section~\ref{S-app:benchmark_hex_square}. Sign-ledger and orientation conventions are recorded in
the supplement, Section~\ref{S-app:sign_ledger}.

%====================================================================
\section{Numerical checks}\label{sec:numerical_diagnostics}
%====================================================================

The numerical material is collected in the online supplement, Section~\ref{S-sec:numerics}, with
reproducibility details in \Cref{S-app:reproducibility_appendix,S-app:minimal_reproduction_script}.
The computations are not proof inputs; they document finite-matrix behavior for the mechanisms
proved above. We retain two numerical figures in the main paper that bear directly on the
headline results: the intrinsic Poincar\'e--Lefschetz comparison
(Figure~\ref{fig:pl_boundary_comparison}) and the finite-\(M\) decoupling log-log scaling
(Figure~\ref{fig:finiteM_loglog}).

\medskip\noindent\textbf{Intrinsic boundary comparison.}
Figure~\ref{fig:pl_boundary_comparison} compares the intrinsic Poincar\'e--Lefschetz boundary
Hamiltonian \(H^{\mathrm{PL}}_+(\K)\), the doubled original-side compression \(H^\partial_+(\K)\),
and the finite-\(M\) jammed model on a disk-like sample. The collar correction
\(R_{\partial,+}=H^\partial_+(\K)-H^{\mathrm{PL}}_+(\K)\) is not norm-small, but it is
confined to the boundary collar predicted by
Theorem~\ref{thm:relative_boundary_comparison}, and the jammed low-energy spectrum tracks the
compressed boundary model in line with the Schur-complement bounds of
Theorem~\ref{thm:resolvent} and Theorem~\ref{thm:first_order_eff}. The numerical refinement scaling of
\(R_{\partial,+}\) -- recorded in the supplement -- is consistent with the boundary-rank
interpretation in Lemma~\ref{lem:collar_rank_counts}.

\begin{figure}[H]
\centering
\safeincludegraphics[width=0.92\linewidth]{fig_pl_boundary_comparison.png}
\caption{Central spectra of \(H^{\mathrm{PL}}_+(\K)\), \(H^\partial_+(\K)\), and
\(H_{+,M}\) on the disk-like example, together with the measured support distribution of the
collar correction \(R_{\partial,+}\). The correction is not norm-small, but it is boundary-collar
supported, as predicted by Theorem~\ref{thm:relative_boundary_comparison}.}
\label{fig:pl_boundary_comparison}
\end{figure}

\medskip\noindent\textbf{Finite-\(M\) log-log scalings.}
Figure~\ref{fig:finiteM_loglog} reports the finite-\(M\) rate diagnostics. The compressed
resolvent error follows the \(M^{-1}\) bound of Theorem~\ref{thm:resolvent}; the first-order
resolvent correction, squared copy leakage, and first-order eigenvalue localization follow the
\(M^{-2}\) predictions of \Cref{thm:first_order_eff,thm:eigvec_leak} and the first-order
eigenvalue localization theorem stated in the supplement,
\Cref{S-app:proof:thm:first_order_spectral_localization}.

\begin{figure}[H]
\centering
\begin{subfigure}{0.49\linewidth}
  \centering
  \safeincludegraphics[width=\linewidth]{fig_resolvent_loglog.png}
  \caption{Compressed resolvent.}
\end{subfigure}\hfill
\begin{subfigure}{0.49\linewidth}
  \centering
  \safeincludegraphics[width=\linewidth]{fig_first_order_resolvent_loglog.png}
  \caption{First-order resolvent.}
\end{subfigure}

\vspace{0.5em}
\begin{subfigure}{0.49\linewidth}
  \centering
  \safeincludegraphics[width=\linewidth]{fig_leakage_loglog.png}
  \caption{Copy leakage.}
\end{subfigure}\hfill
\begin{subfigure}{0.49\linewidth}
  \centering
  \safeincludegraphics[width=\linewidth]{fig_first_order_localization_loglog.png}
  \caption{First-order eigenvalue localization.}
\end{subfigure}
\caption{Log-log finite-\(M\) checks. The compressed resolvent slope tracks the \(M^{-1}\) rate
of Theorem~\ref{thm:resolvent}; the first-order resolvent, squared copy leakage, and first-order
eigenvalue localization track the \(M^{-2}\) rates of \Cref{thm:first_order_eff,thm:eigvec_leak}
and the first-order eigenvalue localization theorem stated in the supplement
(\Cref{S-app:proof:thm:first_order_spectral_localization}). Fitted slopes and additional rate
diagnostics are reported in the supplement, Section~\ref{S-sec:numerics}.}
\label{fig:finiteM_loglog}
\end{figure}

The remaining numerical checks -- boundary transport heatmaps, chirality flip, multiboundary
component drifts, flux insertion, collar/contour/refinement diagnostics, weak-disorder
boundary transport, and the nonconstant positive copy-side confinement -- are all collected in
the supplement, Section~\ref{S-sec:numerics}. A symbolic verification appendix for the local
crossing calculation supporting Theorem~\ref{thm:periodic_triangular_chern} is given in the
supplement, Section~\ref{S-app:symbolic_verification}.

\section{Conclusion}

We constructed a boundary model for the universal Chern Hamiltonians by replacing a triangulated
surface with boundary by its closed double and adding a large positive potential on the reflected
side. The relative fundamental class of \((\K,\bd\K)\) gave an intrinsic Poincar\'e--Lefschetz
boundary Hamiltonian on the original triangulation. Compressing the doubled Hamiltonian to the
original side recovered this relative Hamiltonian up to a boundary-collar operator, and the
jammed Hamiltonian converged to the compressed boundary operator in the large-jam limit by
quantitative Schur-complement estimates.

\paragraph{Three buckets of results.} For an external reader auditing this paper, we partition
the results into three buckets according to which external inputs each one uses. We are explicit
because the headline integer winding has different dependency status in different cases, and a
single label could obscure that.

\begin{itemize}[leftmargin=2em]
\item[(A)] \emph{Fully internal results (no external theorem).} The finite-dimensional
boundary theory (intrinsic Poincar\'e--Lefschetz Hamiltonian, collar comparison, and
Schur-complement estimates of \Cref{thm:relative_boundary_comparison,thm:resolvent,thm:first_order_eff,thm:proj_general,thm:spectral_localization,thm:eigvec_leak,thm:inertia_reduction});
the internal bulk spectral equivalence
(\Cref{thm:internal_bulk_spectral_equivalence}); the periodic triangular benchmark, namely the
exact-determinant gap (\Cref{thm:periodic_triangular_bulk}), the isolated-crossing Chern
calculation (\Cref{thm:periodic_triangular_chern}), and the integer Toeplitz winding in the
straight periodic case (\Cref{cor:toeplitz_winding,cor:periodic_straight_edge_winding}). These
results are proved in the paper without invoking any external theorem and without invoking the
Ewert--Meyer Mayer--Vietoris boundary identification.
\item[(B)] \emph{Internal modulo the Higson--Prodan closed-bulk theorem.} The Roe-algebra
\(K_1\)-class identity \(\Indint(J_{\pm,M}^\infty;T^+,T^-)
=\partial_Y(E_+(\pi(P_\pm)))\in K_1(\mathcal I)\) of \Cref{thm:pathb_roe_winding} relies on
Higson--Prodan's closed-bulk gap and Chern theorem
\cite{HigsonProdanUniversalPRL,HigsonProdanUniversalArxiv}, imported as
\Cref{thm:hp_imported} of Section~\ref{sec:pathb_internalization}. Apart from this single
external input, every step of \Cref{thm:pathb_roe_winding} is proved internally; the
hypothesis check that doubled-and-jammed families satisfy HP's hypotheses is
\Cref{lem:hp_hypothesis_check}.
\item[(C)] \emph{Conditional on an external integer pairing.} The general-partition integer
winding \(\pm1\) statement (\Cref{cor:conditional_winding}) is conditional on the existence of
an integer-valued Chern-jump-compatible pairing between \(K_0(\mathcal A/\mathcal I)\) and
\(K_1(\mathcal I)\). Such a pairing exists internally for the periodic Toeplitz case (already
in bucket (A)); for general partitions the construction of the integer pairing is not carried
out in this paper. This is the only result whose integer value depends on an external
construction beyond the HP closed-bulk theorem, and we flag it explicitly as conditional.
\end{itemize}

\noindent
In particular: the integer winding \(\pm 1\) for the straight periodic triangular interface is
in bucket (A) (no external theorem at all); the \(K_1\)-class identity for general HP-covered
boundary families is in bucket (B); and the integer evaluation of that \(K_1\)-class outside
the periodic Toeplitz case is in bucket (C). No claim in this paper is named in a way that
elides these distinctions; \Cref{cor:toeplitz_winding} carries the qualifier ``in the Toeplitz
periodic case'' and \Cref{cor:conditional_winding} carries the qualifier ``conditional integer
winding for general partitions''.

Beyond the HP-covered and periodic classes treated here, one possible extension is to prove
uniform Fermi gaps for additional closed doubled refinement families. A second direction is to
construct the integer Chern-jump-compatible pairing of bucket (C) for natural classes of
non-Toeplitz partitions, which would promote \Cref{cor:conditional_winding} from bucket (C) to
bucket (B). A third direction is to extend the internal boundary-map identification to curved
or componentwise pairings that separate the contributions of individual boundary components.

\paragraph{Publishable-today status (iter 17 framing).}
Reframing the bucket (A)/(B)/(C) split by dependency on the only cited closed-bulk theorem,
the manuscript naturally separates into two parts according to whether the result still relies
on the Higson--Prodan \cite{HigsonProdanUniversalPRL,HigsonProdanUniversalArxiv} closed-bulk
gap and Chern theorem or whether it is proved internally inside the present paper. We make the
carve-out explicit because it changes the submission framing.
\begin{itemize}[leftmargin=2em]
\item \textbf{Part~A (publishable today, fully HP-internal).} The finite-dimensional
construction (\Cref{sec:setup,sec:decoupling}); the relative Poincar\'e--Lefschetz boundary
Hamiltonian and the doubling-and-jamming setup; the Schur-complement estimates
\Cref{thm:relative_boundary_comparison,thm:resolvent,thm:first_order_eff,thm:proj_general,thm:spectral_localization,thm:eigvec_leak,thm:inertia_reduction};
the periodic flat-torus / hex / square \emph{internal} closed-bulk gap and Chern theorems
of overnight iterations 1--6
(\Cref{S-thm:internal_torus_gap,S-thm:internal_torus_chern,S-thm:internal_hex_gap,S-thm:internal_square_gap,S-thm:internal_hex_chern,S-thm:internal_square_chern});
the iter-11 internal Combes--Thomas decay
\Cref{S-thm:internal_combes_thomas} (conditional only on the abstract gap hypothesis, not on
HP); the periodic triangular Bloch certificate
(\Cref{sec:periodicbulk}, \Cref{thm:periodic_triangular_bulk,thm:periodic_triangular_chern});
the integer Toeplitz winding in the straight periodic case
(\Cref{cor:toeplitz_winding,cor:periodic_straight_edge_winding}); the five topological
diagrams; the numerical experiments. This part is fully internal to the present paper and
does not import the HP closed-bulk theorem for any periodic / flat-torus / hex / square
benchmark. It corresponds to bucket~(A) above plus the iter-1--11 internal benchmark
HP-replacement results and is a substantial publishable contribution on its own.

\item \textbf{Part~B (still conditional on HP, future research).} The global Roe-algebra
interface-winding theorem \Cref{thm:pathb_roe_winding} for arbitrary bounded-geometry
families remains conditional on the Higson--Prodan closed-bulk gap of
\Cref{thm:hp_imported}(i) and the Higson--Prodan closed-bulk Chern integer of
\Cref{thm:hp_imported}(ii). The general non-periodic, non-Bravais case is therefore Part~B.
The iter-12 inventory \Cref{S-app:hp_inventory_iter12} confirms that the bounded-geometry
closed-bulk gap of \Cref{thm:hp_imported}(i) is the single remaining research bottleneck;
the iter-13 completeness verification \Cref{S-app:completeness_iter13} confirms that the
seven internal HP-replacement results of Part~A do \emph{not} cover this bottleneck. Part~B
is the natural target of a future research-medium iteration; it is the only part of the
manuscript whose status depends on whether HP is treated as a citation or is re-derived
internally. The iter-23 supplement subsection
\Cref{S-app:delaunay_first_step} records a partial first step in this
direction: it defines the quasi-uniform Delaunay class
\(\mathcal D(\theta_0,L)\) (a strict sub-class of the bounded-geometry
class, parametrized by an explicit minimum-angle bound \(\theta_0>0\)
and an edge-length ratio bound \(L\ge 1\)), proves
(\Cref{S-lem:delaunay_bg_parameters}) that every triangulation in this
class is bounded-geometry with constants depending only on
\((\theta_0,L)\), and states as a conjecture with proof outline
(\Cref{S-thm:delaunay_gap_stability_conjecture}) the relative stability
of the closed-bulk gap under Delaunay-class perturbations; the
absolute open question (\Cref{S-rem:delaunay_open_question}) of
whether at least one non-periodic quasi-uniform Delaunay triangulation
can be shown internally to have a positive Fermi gap remains genuinely
open. The iter-23 deliverable is the proved bounded-geometry parameter
lemma plus the framework for the conjecture; the conjecture itself is
not proved. The iter-24 supplement subsection
\Cref{S-app:cap_product_coherence_iter24} records a partial attempt at
the first of the two missing ingredients identified in the iter-23
``What is missing in the outline'' paragraph (the cap-product coherence
of \(S_K\) under quasi-isometric pullback). The iter-24 attempt proves
the simplicial-isomorphism / translation sub-case rigorously
(\Cref{S-lem:translation_coherence_iter24}: for two simplicially
isomorphic Delaunay triangulations, the universal Hamiltonians are
exactly unitarily conjugate, hence the Fermi gaps are identical) and a
vertex-perturbation extension (\Cref{S-lem:vertex_perturbation_coherence_iter24}:
the cap-product operator \(S_K\) is purely combinatorial and is
unchanged by vertex-position perturbations within a fixed isotopy
class); the general quasi-isometric case is recorded as a refined
conjecture (\Cref{S-thm:qi_coherence_general_iter24}) with an explicit
support-and-norm decomposition. The second iter-23 missing ingredient
(the isometric isomorphism between Hilbert spaces of different
dimension) is not addressed in iter 24 and remains open. The iter-25
supplement subsection
\Cref{S-app:hilbert_iso_diff_dim_iter25} records a partial attempt
at the \emph{second} of the two missing ingredients identified in
the iter-23 ``What is missing in the outline'' paragraph (the
isometric isomorphism between Hilbert spaces of different dimension).
The iter-25 attempt defines an explicit \emph{inclusion partial
isometry} \(V_{K\hookrightarrow K'}:\Hil(K)\to\Hil(K')\) for
sub-complex pairs \(K\subseteq K'\)
(\Cref{S-def:inclusion_partial_isometry_iter25}), proves a
boundary-supported coherence lemma
(\Cref{S-lem:inclusion_coherence_iter25}:
\(V^*\,S_{K'}\,V=S_K+\Delta_V\) with
\(\mathrm{supp}(\Delta_V)\subseteq\partial_{K'}K\) and operator-norm
bound \(\|\Delta_V\|\le 2\Delta_*(\theta_0)\) where
\(\Delta_*(\theta_0)\) is the bounded-geometry vertex-link constant
of \Cref{S-lem:delaunay_bg_parameters}), and records the general
common-refinement case as a refined conjecture
(\Cref{S-thm:diff_dim_coherence_iter25}) with an explicit
support-and-norm decomposition. The iter-25 remark
\Cref{S-rem:comparison_iter24_iter25} explicitly records how the
iter-24 and iter-25 attempts combine through an
inclusion + simplicial + reverse-inclusion decomposition to give a
conditional resolution of the iter-23 stability conjecture, modulo
the open question of common-refinement existence
(\Cref{S-rem:inclusion_vs_refinement_iter25}). The iter-27
supplement subsection
\Cref{S-app:single_flip_closed_form_iter27} closes sub-case~(b)
of the iter-24 conjectural general QI theorem
\Cref{S-thm:qi_coherence_general_iter24} (the single-flip local
edit case) by an explicit closed-form computation of the
cap-product matrix elements on the eleven-simplex local Hilbert
space of a Delaunay flip quadrilateral, together with a Schur-test
operator-norm bound
\(\|R_\phi^{\mathrm{single}}\|\le 4\)
(\Cref{S-thm:single_flip_schur_iter27}), with the explicit
constant \(C_{\mathrm{flip}}=4\) independent of
\((\theta_0,L)\). The iter-28 supplement subsection
\Cref{S-app:multi_flip_schur_iter28} closes sub-case~(c) of the
iter-24 conjectural theorem (the multi-flip assembly) by a
telescoping decomposition of the cumulative correction operator
into \(N_{\mathrm{flip}}\) single-flip pieces and a
triangle-inequality bound
\(\|R_\phi\|\le C_{\mathrm{flip}}\,N_{\mathrm{flip}}=4N_{\mathrm{flip}}\)
(\Cref{S-thm:multi_flip_schur_iter28}), with an edit-graph
bounded-incidence lemma
(\Cref{S-lem:bounded_edit_incidence_iter28}) supporting a sharper
Schur-test refinement in the star-disjoint case. Combining
sub-cases~(a) (iter 24: \Cref{S-lem:translation_coherence_iter24}),
(b) (iter 27:
\Cref{S-thm:single_flip_schur_iter27}), and~(c) (iter 28:
\Cref{S-thm:multi_flip_schur_iter28}), the iter-24 conjectural
theorem \Cref{S-thm:qi_coherence_general_iter24} is upgraded
from \textbf{conjectural} to \textbf{proved internally} at iter 28
with the explicit constant
\(C_*(\theta_0,L)=C_{\mathrm{flip}}=4\) (independent of
\((\theta_0,L)\)), strictly tighter than the iter-24 provisional
\(8\,\Delta_*(\theta_0)\,L_*(\theta_0)\). Ingredient~(i) of the
iter-23 missing-ingredients list (the cap-product coherence
under quasi-isometric pullback) is therefore now closed in full
generality. The iter-23 stability conjecture
\Cref{S-thm:delaunay_gap_stability_conjecture} nonetheless
remains conjectural at iter 28 because ingredient~(ii) (the
iter-25 general different-dimension coherence theorem
\Cref{S-thm:diff_dim_coherence_iter25}) remains conjectural.
The iter-29 supplement subsection
\Cref{S-app:common_refinement_iter29} proposes an explicit
candidate construction for the common refinement \(K_*\)
required by iter-25 sub-case~(c): the \emph{union-Delaunay
refinement} \(K_*=K_*(K,K')\)
(\Cref{S-def:union_delaunay_iter29}) is the Delaunay
triangulation of the union vertex set \(V_K\cup V_{K'}\).
Vertex inclusion \(V_K\subseteq V_{K_*}\),
\(V_{K'}\subseteq V_{K_*}\) is proved unconditionally
(\Cref{S-lem:union_delaunay_vertices_iter29}), and a parameter
relaxation
\((\theta_0,L)\to(\theta_0',L')\) with explicit bounds
depending on the minimal inter-vertex distance
\(d_{\min}(V_K,V_{K'})\) and the maximum edge length
\(\overline{\ell}_{\max}\) is proved unconditionally
(\Cref{S-lem:union_delaunay_parameters_iter29}). The
union-Delaunay refinement provides \emph{vertex} inclusion but
\emph{not} the simplicial sub-complex inclusion required by
\Cref{S-def:inclusion_partial_isometry_iter25}; a residual
sub-complex-refinement existence sub-conjecture
(\Cref{S-rem:subcomplex_refinement_obstruction_iter29}) is
recorded honestly and is plausible by barycentric subdivision
followed by Delaunay edge-flip repair, but is not proved at
iter 29. Conditional on this sub-conjecture, the three-step
composition theorem
\Cref{S-thm:three_step_composition_iter29} closes iter-25
sub-case~(c) with explicit constant
\(C_\diamond(\widetilde\theta_0,\widetilde L)
=6\,\Delta_*(\widetilde\theta_0)^2\,L_*(\widetilde\theta_0)\)
in the relaxed parameter class. Combining the iter-28
unconditional closure of ingredient~(i) with the iter-29
conditional closure of ingredient~(ii), the iter-23 stability
conjecture
\Cref{S-thm:delaunay_gap_stability_conjecture} is upgraded
from \textbf{conjectural} (iters 23--28) to
\textbf{conditional} (iter 29, modulo
\Cref{S-rem:subcomplex_refinement_obstruction_iter29}), with
explicit Combes--Thomas locality constants inherited from
\Cref{S-lem:delaunay_bg_parameters} applied in the relaxed
parameter class (the precise conditional statement is recorded
in
\Cref{S-rem:iter29_iter23_conditional_iter29}). The iter-30
supplement subsection
\Cref{S-app:subcomplex_refinement_iter30} attempts the natural
weakened-inclusion route to closing the iter-29 sub-conjecture
by replacing the simplicial sub-complex relation
\(K\subseteq K_*\) of
\Cref{S-def:inclusion_partial_isometry_iter25} with a
\emph{refinement-in-the-union sense} relation
\(K\preceq^\cup\widetilde K_*\) of
\Cref{S-def:refinement_union_iter30}, in which every simplex of
\(K\) decomposes as a union of simplices in
\(\widetilde K_*\). The corresponding \emph{block partial
isometry} \(V_K^\cup:\mathcal H(K)\to\mathcal H(\widetilde K_*)\)
of \Cref{S-def:block_partial_isometry_iter30} is rigorously
proved to be isometric
(\Cref{S-lem:block_isometry_iter30}), and an explicit
constrained common refinement (the overlay of \(K\) and
\(K'\)) is constructed
(\Cref{S-lem:overlay_refinement_iter30}) with explicit
maximum-split-factor bound
\(N_{\max}\le 2L_*(\theta_0)^2\). However, the iter-30
\Cref{S-lem:block_capproduct_matrix_iter30} identifies a
concrete renormalization mismatch factor
\(\sqrt{N(\sigma')/N(\sigma)}\) on the cap-product matrix
elements of \(V^*\,S_{\widetilde K_*}\,V\) for incident pairs
\(\sigma'\subset\sigma\) in \(K\): when \(\sigma\) and
\(\sigma'\) have different split factors, the matrix element of
\((V_K^\cup)^*\,S_{\widetilde K_*}\,V_K^\cup\) differs from that
of \(S_K\) by a multiplicative factor
\(\sqrt{N(\sigma')/N(\sigma)}\), \emph{not} by an additive
boundary-supported correction. The induced coherence
correction \(\Delta_V^\cup=(V_K^\cup)^*\,S_{\widetilde K_*}\,V_K^\cup-S_K\)
is therefore \emph{bulk-supported}, not boundary-supported, with
global operator-norm bound
\(\|\Delta_V^\cup\|\le 2\sqrt{2}\,\Delta_*(\theta_0')\,L_*(\theta_0)\)
(\Cref{S-thm:block_obstruction_iter30}). The bulk support
\emph{prevents} direct application of the iter-25
inclusion-coherence lemma
\Cref{S-lem:inclusion_coherence_iter25} in the iter-29
three-step composition theorem
\Cref{S-thm:three_step_composition_iter29}. Iter 30 is therefore
a \emph{negative-result} iteration: it identifies a concrete and
honest obstruction to the natural weakened-inclusion approach,
illustrated explicitly on the flat torus
(\Cref{S-rem:flat_torus_block_obstruction_iter30}). The iter-23
stability cascade status is \emph{unchanged from iter 29}: the
iter-23 stability conjecture
\Cref{S-thm:delaunay_gap_stability_conjecture} remains
\textbf{conditional} (modulo the iter-29 sub-conjecture
\Cref{S-rem:subcomplex_refinement_obstruction_iter29}), \emph{not}
upgraded to unconditional at iter 30; the honest status is
recorded in
\Cref{S-rem:iter30_cascade_honest_iter30}.
\end{itemize}

The full Part~A vs.\ Part~B carve-out, with each theorem placed in the appropriate column
and a recommended submission strategy, is given in \Cref{S-app:publishable_today_iter17}.
Concretely, Part~A is publishable as a standalone substantial contribution to a serious
mathematical-physics journal: a finite-dimensional doubling-and-jamming framework with
periodic-benchmark internal closed-bulk gap and Chern theorems, an internal Combes--Thomas
decay lemma, and integer Toeplitz windings on the periodic triangular interface. Part~B is
honestly framed as a conditional theorem whose closed-bulk inputs remain external HP and
whose bounded-geometry case is identified as the future research push.

\section*{Code and data availability}\label{sec:code_availability}
The companion numerical code and the online supplement are available at the project repository,
which is hosted on GitHub:
\href{https://github.com/yspennstate/Doubling-and-Jamming-A-Universal-Chern-Hamiltonian-on-Triangulated-Surfaces-with-Boundary}{\texttt{https://github.com/yspennstate/Doubling-and-Jamming-A-Universal-Chern-Hamiltonian-on-Triangulated-Surfaces-with-Boundary}}.
The online supplement is also distributed through that same GitHub repository. The
periodic triangular Bloch certificate script
(\texttt{periodic\_triangular\_bulk\_certificate.py}) and the companion numerical pipeline
that generates the figures and tables of the supplement are included.

\section*{Acknowledgments}
The research presented in this paper was supported by the European Research Council (ERC) under the European Union's Horizon Europe research and innovation programme (grant agreement No.~101041711), by the Simons Foundation (as part of the Collaboration on the Mathematical and Scientific Foundations of Deep Learning), by the Israel Science Foundation (grant number 2258/19), by the Israel Science Foundation (ISF Grant~4101/25), and by the U.S. National Science Foundation (NSF Grant~OISE-2401227).

\section*{Declaration of generative AI and AI-assisted technologies in the manuscript preparation process}
During the preparation of this work the author used generative AI tools to accelerate drafting
and revision. All mathematical claims, proofs, and bibliographic information were subsequently
reviewed and edited by the author, who takes full responsibility for the content of the manuscript.

\FloatBarrier

\begin{thebibliography}{99}

\bibitem{AsbothOroszlanyPalyi2016}
J.~K.~Asboth, L.~Oroszlany, and A.~Palyi,
\emph{A Short Course on Topological Insulators: Band Structure and Edge States in One and Two
Dimensions},
Lecture Notes in Physics \textbf{919}, Springer, 2016.

\bibitem{BellissardVanElstSchulzBaldes1994}
J.~Bellissard, A.~van Elst, and H.~Schulz-Baldes,
\emph{The noncommutative geometry of the quantum Hall effect},
J. Math. Phys. \textbf{35} (1994), 5373--5451.

\bibitem{Haldane1988}
F.~D.~M.~Haldane,
\emph{Model for a Quantum Hall Effect without Landau Levels: Condensed-Matter Realization of the Parity Anomaly},
Phys. Rev. Lett. \textbf{61} (1988), 2015--2018.

\bibitem{SchnyderRyuFurusakiLudwig}
A.~P.~Schnyder, S.~Ryu, A.~Furusaki, and A.~W.~W.~Ludwig,
\emph{Classification of topological insulators and superconductors in three spatial dimensions},
Phys. Rev. B \textbf{78} (2008), 195125.

\bibitem{KitaevPeriodic}
A.~Kitaev,
\emph{Periodic table for topological insulators and superconductors},
AIP Conf. Proc. \textbf{1134} (2009), 22--30.

\bibitem{QiZhangRMP}
X.-L.~Qi and S.-C.~Zhang,
\emph{Topological insulators and superconductors},
Rev. Mod. Phys. \textbf{83} (2011), 1057--1110.

\bibitem{FukuiHatsugaiSuzuki}
T.~Fukui, Y.~Hatsugai, and H.~Suzuki,
\emph{Chern numbers in discretized Brillouin zone: efficient method of computing (spin) Hall conductances},
J. Phys. Soc. Jpn. \textbf{74} (2005), 1674--1677.

\bibitem{HigsonProdanUniversalPRL}
N.~Higson and E.~Prodan,
\emph{Universal Chern Model on Arbitrary Triangulations},
Phys. Rev. Lett. \textbf{136} (2026), no.~9, 096602,
doi:10.1103/66x5-dvqq.

\bibitem{HigsonProdanUniversalArxiv}
N.~Higson and E.~Prodan,
\emph{A Universal Chern Model on Arbitrary Triangulations},
arXiv:2510.20862, doi:10.48550/arXiv.2510.20862.

\bibitem{HigsonRoeI}
N.~Higson and J.~Roe,
\emph{Mapping Surgery to Analysis I: Analytic Signatures},
K-Theory \textbf{33} (2004), 277--299.

\bibitem{HigsonRoeII}
N.~Higson and J.~Roe,
\emph{Mapping Surgery to Analysis II: Geometric Signatures},
K-Theory \textbf{33} (2004), 301--324.

\bibitem{HigsonRoeIII}
N.~Higson and J.~Roe,
\emph{Mapping Surgery to Analysis III: Exact Sequences},
K-Theory \textbf{33} (2004), 325--346.

\bibitem{WallSurgery}
C.~T.~C.~Wall,
\emph{Surgery on Compact Manifolds},
AMS, 1999.

\bibitem{MunkresElements}
J.~R.~Munkres,
\emph{Elements of Algebraic Topology},
Addison-Wesley, 1984.

\bibitem{Hatcher}
A.~Hatcher,
\emph{Algebraic Topology},
Cambridge University Press, 2002.

\bibitem{KatoPerturbation}
T.~Kato,
\emph{Perturbation Theory for Linear Operators},
Classics in Mathematics, Springer, 1995.

\bibitem{BhatiaMatrixAnalysis}
R.~Bhatia,
\emph{Matrix Analysis},
Graduate Texts in Mathematics, vol.~169, Springer, 1997.

\bibitem{AtiyahKTheory}
M.~F.~Atiyah,
\emph{K-Theory},
W.~A.~Benjamin, 1967.

\bibitem{BlackadarKTheory}
B.~Blackadar,
\emph{K-Theory for Operator Algebras},
2nd ed., Mathematical Sciences Research Institute Publications, vol.~5,
Cambridge University Press, 1998.

\bibitem{BDF1977}
L.~G.~Brown, R.~G.~Douglas, and P.~A.~Fillmore,
\emph{Extensions of $C^*$-algebras and $K$-homology},
Ann. of Math. (2) \textbf{105} (1977), no.~2, 265--324.

\bibitem{ProdanSchulzBaldesBook}
E.~Prodan and H.~Schulz-Baldes,
\emph{Bulk and Boundary Invariants for Complex Topological Insulators: From K-Theory to Physics},
Springer, 2016.

\bibitem{Mitchell2018}
N.~P.~Mitchell, L.~M.~Nash, D.~Hexner, A.~Turner, and W.~T.~M.~Irvine,
\emph{Amorphous topological insulators constructed from random point sets},
Nat. Phys. \textbf{14} (2018), 380--385.

\bibitem{BourneProdan2018}
C.~Bourne and E.~Prodan,
\emph{Non-commutative Chern numbers for generic aperiodic discrete systems},
J. Phys. A: Math. Theor. \textbf{51} (2018), 235202.

\bibitem{BourneMesland2019}
C.~Bourne and B.~Mesland,
\emph{Index theory and topological phases of aperiodic lattices},
Ann. Henri Poincar\'e \textbf{20} (2019), 1969--2038.

\bibitem{ElbauGraf2002}
P.~Elbau and G.~M.~Graf,
\emph{Equality of Bulk and Edge Hall Conductance Revisited},
Comm. Math. Phys. \textbf{229} (2002), 415--432.

\bibitem{ElgartGrafSchenker2005}
A.~Elgart, G.~M.~Graf, and J.~H.~Schenker,
\emph{Equality of the Bulk and Edge Hall Conductances in a Mobility Gap},
Comm. Math. Phys. \textbf{259} (2005), 185--221.

\bibitem{GrafPorta2013}
G.~M.~Graf and M.~Porta,
\emph{Bulk-Edge Correspondence for Two-Dimensional Topological Insulators},
Comm. Math. Phys. \textbf{324} (2013), 851--895.

\bibitem{KubotaControlled}
Y.~Kubota,
\emph{Controlled topological phases and bulk-edge correspondence},
Comm. Math. Phys. \textbf{349} (2017), 493--525.

\bibitem{EwertMeyer2019}
E.~E.~Ewert and R.~Meyer,
\emph{Coarse geometry and topological phases},
Comm. Math. Phys. \textbf{366} (2019), 1069--1098,
doi:10.1007/s00220-019-03303-z.

\bibitem{BourneKellendonkRennie}
C.~Bourne, J.~Kellendonk, and A.~Rennie,
\emph{The $K$-theoretic bulk-edge correspondence for topological insulators},
Ann. Henri Poincar\'e \textbf{18} (2017), 1833--1866.

\bibitem{AlldridgeMaxZirnbauer2020}
A.~Alldridge, C.~Max, and M.~R.~Zirnbauer,
\emph{Bulk-Boundary Correspondence for Disordered Free-Fermion Topological Phases},
Comm. Math. Phys. \textbf{377} (2020), 1761--1821.

\bibitem{LudewigThiangEdge}
M.~Ludewig and G.~C.~Thiang,
\emph{Cobordism invariance of topological edge-following states},
Adv. Theor. Math. Phys. \textbf{26} (2022), no.~3, 673--710.

\bibitem{DrouotZhuEdgeSpectrum}
A.~Drouot and X.~Zhu,
\emph{Topological edge spectrum along curved interfaces},
Int. Math. Res. Not. IMRN \textbf{2024} (2024), no.~22, 13870--13889.

\bibitem{DrouotZhuBEC}
A.~Drouot and X.~Zhu,
\emph{The bulk-edge correspondence for curved interfaces},
arXiv:2408.07950.

\bibitem{RoeLectures}
J.~Roe,
\emph{Lectures on Coarse Geometry},
University Lecture Series \textbf{31}, American Mathematical Society, 2003.

\bibitem{CombesThomas1973}
J.~M.~Combes and L.~Thomas,
\emph{Asymptotic behaviour of eigenfunctions for multiparticle Schr\"odinger operators},
Comm. Math. Phys. \textbf{34} (1973), 251--270.

\bibitem{AizenmanWarzel2015}
M.~Aizenman and S.~Warzel,
\emph{Random Operators: Disorder Effects on Quantum Spectra and Dynamics},
Graduate Studies in Mathematics, vol.~168, American Mathematical Society, 2015.

\bibitem{Cheeger1970}
J.~Cheeger,
\emph{A lower bound for the smallest eigenvalue of the Laplacian},
in \emph{Problems in Analysis (Papers Dedicated to Salomon Bochner, 1969)},
Princeton University Press, 1970, pp.~195--199.

\bibitem{Chung1997}
F.~R.~K.~Chung,
\emph{Spectral Graph Theory},
CBMS Regional Conference Series in Mathematics, vol.~92,
American Mathematical Society, 1997.

\bibitem{ChavelEigenvalues}
I.~Chavel,
\emph{Eigenvalues in Riemannian Geometry},
Pure and Applied Mathematics, vol.~115, Academic Press, 1984.

\bibitem{Petersen2016}
P.~Petersen,
\emph{Riemannian Geometry},
3rd ed., Graduate Texts in Mathematics, vol.~171, Springer, 2016.

\bibitem{LottVillani2009}
J.~Lott and C.~Villani,
\emph{Ricci curvature for metric-measure spaces via optimal transport},
Ann. of Math. (2) \textbf{169} (2009), no.~3, 903--991.

\bibitem{AdamsFournier}
R.~A.~Adams and J.~J.~F.~Fournier,
\emph{Sobolev Spaces},
2nd ed., Pure and Applied Mathematics, vol.~140, Academic Press, 2003.

\bibitem{Strang1973}
G.~Strang and G.~J.~Fix,
\emph{An Analysis of the Finite Element Method},
Prentice-Hall, 1973.

\end{thebibliography}

\end{document}
